% concatenated from main.tex
\documentclass[12pt]{article}
\usepackage[utf8]{inputenc}

\usepackage[left=1cm, top=1cm, right=1cm, bottom=1cm, includeheadfoot, headheight=0pt, a4paper]{geometry}
\usepackage{fancyhdr}
\usepackage{lastpage}
\usepackage{amsmath}
\usepackage{amssymb}
\usepackage{graphicx}
\usepackage{enumitem}
\usepackage{multicol}
%\usepackage{amsart}
\usepackage{amsmath}
\usepackage{amsfonts}
%\usepackage{polski}
\usepackage[utf8]{inputenc}
\usepackage{amsthm}
\usepackage{tikz-cd}
\usepackage{geometry}
\usepackage{mathtools}
\usepackage{tikz}
\usepackage{comment}
\usepackage{url}
\usetikzlibrary{patterns}
\usepackage{svg}
\usepackage{xcolor}
\usepackage[T1]{fontenc}
\usepackage[utf8]{inputenc}
\usepackage{float}
\usepackage{subcaption}

\newtheorem{thm}{Theorem}[section]
\newtheorem{lemma}[thm]{Lemma}
\newtheorem{df}[thm]{Definition}
\newtheorem{prop}[thm]{Proposition}
\newtheorem{cor}[thm]{Corollary}
\newtheorem{fact}[thm]{Fact}
\newtheorem{conj}[thm]{Conjecture}
\newtheorem{prob}[thm]{Problem}
\theoremstyle{remark}
\newtheorem{remark}[thm]{Remark}
\newtheorem{ex}[thm]{Example}
\renewcommand{\thefootnote}{\fnsymbol{footnote}}

\newcommand\blfootnote[1]{%
  \begingroup
  \renewcommand\thefootnote{}\footnote{#1}%
  \addtocounter{footnote}{-1}%
  \endgroup
}



\pagestyle{fancy}

\fancyhf{}

\cfoot{
$\thepage$
}

\renewcommand{\headrulewidth}{0pt}



\title{The type and cardinality of minimal presentations of \\ numerical semigroups with embedding dimension four }

\author{Kazimierz Chomicz}
\date{September 2026}
\begin{document}
%%%%%%%%%%%%%%%%%%%%%%%%%%%%%%%%%%%%%%%%%%%%%%%%%%%%%%%%%%%%%%%%%%%%%%%%%%%

\maketitle



\abstract{ 
For a numerical semigroup $S$ with embedding dimension four, we study the relationship between its type $t(S)$ and the cardinality of its minimal presentations $\eta(S)$. Using an approach based on the geometry of the Ap{\'e}ry set, we prove that $4t(S) + 5 \geq \eta(S) \geq t(S)-11$. This resolves a problem of Moscariello and Sammartano asking whether $t(S)$ is bounded by a function of $\eta(S)$, and improves the previously known bound $9t(S) + 4 \geq \eta(S)$ due to Bresinsky.
}




\section{Introduction}

\blfootnote{\text{Keywords: numerical semigroup, embedding dimension, Ap{\'e}ry set, type, minimal presentation.}}
\blfootnote{\text{2020 Mathematics Subject Classification: 20M13, 20M14}}

Let $\mathbb{N}$ denote the set of nonnegative integers. For $k> 0$ and $d_0 ,d_1, \dots, d_k \in \mathbb{N}$ such that $\gcd(d_0,d_1, \dots, d_k) = 1$, the set
$$ S = \langle d_0,d_1, \dots, d_k \rangle  = \{ \lambda_0 d_0 + \lambda_1 d_1 + \dots + \lambda_k d_k \mid \lambda_i \in \mathbb{N}\},$$
is called a \emph{numerical semigroup} generated by $\{d_0,d_1,\dots,d_k\}$. Equivalently, $S$ is a submonoid of $(\mathbb{N},+)$ such that $\mathbb{N} \setminus S$ is finite (see \cite[Lemma 2.1]{Rosales2009}). If no proper subset of $\{d_0,d_1,\dots,d_k\}$ generates $S$, then we say that $S$ is minimally generated by $\{d_0,d_1,\dots,d_k\}$. The cardinality of the minimal set of generators of $S$ is called the \emph{embedding dimension} of $S$ and is denoted $e(S)$. 

The largest integer that is not in $S$ is called the \emph{Frobenius number} and is denoted $F(S)$. An integer $a \notin S$ is called a \emph{pseudo-Frobenius number} of $S$ if $a + s \in S$ for all $s \in S \setminus \{0\}$. The set of all pseudo-Frobenius numbers of $S$ is denoted $PF(S)$. The cardinality of $PF(S)$ is called the \emph{type} of $S$ and is denoted $t(S)$. 

If $S$ is minimally generated by $\{d_0,d_1,\dots,d_k\}$, then we have a homomorphism $\varphi : \mathbb{N}^{k+1} \to S$ defined by $\varphi(z_0, \dots, z_k) = z_0d_0 + \dots + z_kd_k$. 
The \emph{kernel congruence} of $\varphi$ is $\ker(\varphi) = \{ (a,b) \in \mathbb{N}^{k+1} \times \mathbb{N}^{k+1} \mid \varphi(a) = \varphi(b) \}$. 
It is well known that $S$ is isomorphic to the monoid $\mathbb{N}^{k+1} / \ker(\varphi)$. A minimal set of generators of the congruence $\ker(\varphi)$ is called a \emph{minimal presentation} of $S$. Minimal presentations of $S$ are not unique, but their cardinality depends only on $S$ and is denoted by $\eta(S)$. 

Both type and minimal presentations are central to the study of numerical semigroups (see \cite{AssiDAnnaGarciaSanchez2020,BarucciDobbsFontana1997,Froberg1987,MoscarielloSammartano2025,Rosales2009}).
In the study of minimal presentations, the invariant $\eta(S)$ has attracted considerable interest (see \cite{Elmacioglu2024,MoscarielloSammartanominpresen2025,MoscarielloSammartano2025,Rosales2009,Stamate2018}). 

A numerical semigroup $S$ is \emph{symmetric} if for every $x \in \mathbb{Z} \setminus S$ we have $F(S) - x \in S$. This is equivalent to $t(S)=1$ \cite{Kunz1970}. It is well known that if $e(S)=2$, then $t(S) = \eta(S)=1$. The case $e(S)=3$ is completely settled in the following theorem.


\begin{thm}\label{theorem_embd3}
    \textup{\cite{Herzog1970}} Let $S$ be a numerical semigroup with $e(S)=3$. If $S$ is symmetric, then $\eta(S)=2$ and $t(S)=1$. If $S$ is not symmetric, then $\eta(S)=3$ and $t(S)=2$.
\end{thm}

When $e(S) \geq 4$, the situation becomes much more complicated. Bresinsky \cite{Bresinsky1975onprimeideals} found a family of numerical semigroups with $e(S) = 4$ whose $\eta(S)$ can be arbitrarily large, while Backelin \cite[p.\ 75]{Froberg1987} found a family of numerical semigroups with $e(S) = 4$ whose $t(S)$ can be arbitrarily large.

In this paper, we study $t(S)$ and $\eta(S)$ for numerical semigroups with embedding dimension four. In \cite{Bresinsky1988}, Bresinsky established the following bound.

\begin{thm}\label{theoremtypebre}
    \textup{\cite[Theorem 7]{Bresinsky1988}} If $S$ is a numerical semigroup with $e(S) = 4$, then $  9 t(S) + 4 \geq \eta(S)$.
\end{thm}

In \cite{MoscarielloSammartano2025}, the authors, inspired by the above bound, proposed the following problem.

\begin{prob}\label{problemtype}
    \textup{\cite[Problem 25]{MoscarielloSammartano2025}} Let $S$ be a numerical semigroup with $e(S) = 4$. Is $t(S)$ bounded by a function of $\eta(S)$?
\end{prob}

Resolving Problem \ref{problemtype} also settles a general
problem on the boundedness of Betti numbers \cite[Problem 24]{MoscarielloSammartano2025} in the case $e(S)=4$, as noted in the same paper.

The main result of this paper is the following theorem, which settles Problem \ref{problemtype} and improves the bound of Theorem \ref{theoremtypebre}.

\begin{thm}\label{theorem_type_main}
    If $S$ is a numerical semigroup with $e(S) = 4$, then $4 t(S) +5 \geq \eta(S) \geq t(S) - 11$. 
\end{thm}

 

The paper is organized as follows.
In Section \ref{section_geometricprocedure}, we study the geometry of the Ap{\'e}ry set of numerical semigroups with embedding dimension four using $L$-shapes \cite{AGUILOGOST2015}. 
In Section \ref{section_minimal_presentations}, we relate minimal presentations of numerical semigroups with embedding dimension four to an $L$-shape.
In Section \ref{section_type}, we prove Theorem \ref{theorem_type_main} by using results of the previous sections to study $t(S)$ and $\eta(S)$ geometrically.
We conclude the paper with Conjecture \ref{conjecture}.






\section{Geometry of the Ap{\'e}ry set}\label{section_geometricprocedure}



In this section, we study the geometry of the Ap{\'e}ry set of numerical semigroups with $e(S)=4$. Subsection~\ref{subsection_minimal_relations} introduces the notion of minimal relations and uses them to classify numerical semigroups with $e(S)=4$ into primary and secondary semigroups (Lemma~\ref{lemma2cases}) --- a distinction that underpins our approach. Subsection~\ref{subsection_framework} presents the geometric framework for analyzing the Apéry set (Theorem~\ref{theoremprocedure}). Subsection~\ref{subsection_Lshapes} introduces the notion of $L$-shapes and examines them for primary and secondary semigroups (Propositions \ref{prop_primary} and \ref{prop_secondary}).

Let $S$ be minimally generated by $\{d_0,d_1,\dots,d_k\}$. For an element $s \in S$, the set
\begin{equation}
\varphi^{-1}(s) = \{(z_0,\dots,z_k)\in \mathbb{N}^{k+1} \mid z_0d_0 + \dots + z_kd_k = s\}
\end{equation}
is called the \emph{set of factorizations} of $s$.
For $z =(z_0,\dots,z_k) \in \varphi^{-1}(s)$, the \emph{support} of $z$ is 
\begin{equation}
\text{supp}(z) = \{i \in \{0,\dots,k\} \mid z_i > 0\}.
\end{equation}
For an element $a \in S$, the set
\begin{equation}
\textup{Ap}(S,a) = \{ s \in S \mid s-a \notin S\}
\end{equation}
is called the \emph{Ap{\'e}ry set} of $S$ with respect to $a$. Note that $\textup{Ap}(S,a)$ is the set of the smallest nonnegative integers from each residue class modulo $a$ that are in $S$, hence $ |\textup{Ap}(S,a)| = a$. From \cite[Proposition~2.20]{Rosales2009}, we have
\begin{equation}\label{eq_pseudo_frob}
     PF(S) = \{s - a \mid s \in \textup{Ap}(S,a), s + d_i \notin \textup{Ap}(S,a) \text{ for } i = 0,1,\dots,k\}.
\end{equation}
This description of pseudo-Frobenius numbers in terms of the Ap{\'e}ry set will provide a useful framework for our study of the type in Section \ref{section_type}.





\subsection{Minimal relations}\label{subsection_minimal_relations}

Let $S =\langle d_0,d_1,d_2,d_3 \rangle$ be a numerical semigroup with embedding dimension four.

For every $i \in \{0,1,2,3\}$, let $a_{ii}$ be the smallest positive integer such that there exist nonnegative integers $a_{ij},a_{ik},a_{il}$ satisfying
\begin{equation}\label{minimalrelation_eq1}
     a_{ii} d_i = a_{ij} d_j + a_{ik}d_k + a_{il} d_l
\end{equation}
where $\{i,j,k,l\} = \{0,1,2,3\}$. We call \eqref{minimalrelation_eq1} a \emph{minimal relation} of $d_i$.

Fix $m \in \{0,1,2,3\}$. For every $i \in \{0,1,2,3\} \setminus \{m\}$, let $b_{m,ii}$ be the smallest positive integer such that there exist nonnegative integers $b_{m,ij},b_{m,ik},b_{m,il}$ satisfying 
\begin{equation}\label{minimalrelation_eq2}
     b_{m,ii} d_i = b_{m,ij} d_j + b_{m,ik}d_k + b_{m,il} d_l
\end{equation}
where $\{i,j,k,l\} = \{0,1,2,3\}$ and $b_{m,im} > 0$. We call \eqref{minimalrelation_eq2} a \emph{$d_m$-positive minimal relation of $d_i$}.

Note that in the minimal relations $a_{ii}$ is unique; however, each $a_{ij}$ for $i \neq j$ is not necessarily unique. Similarly, in the $d_m$-positive minimal relations $b_{m,ii}$ is unique; however, each $b_{m,ij}$ for $i \neq j$ is not necessarily unique.

Let $\varphi : \mathbb{N}^{4} \to S$ be the factorization morphism of $S$. 

\begin{lemma}\label{lemma_catenary_minrelations}
    For $i \in \{0,1,2,3\}$ and $(\mu_0,\mu_1,\mu_2,\mu_3) \in \varphi^{-1}( a_{ii}d_i)$, we have either $\mu_i = a_{ii}$ and $\mu_j = 0$ for all $j \in \{0,1,2,3\} \setminus \{i\}$, or $\mu_i = 0$.
\end{lemma}

\begin{proof}
     Without loss of generality, suppose $i=0$. We have
     \begin{equation}\label{lemma_catenary2_eq1}
         (a_{00}-\mu_0)d_0 =  \mu_1d_1 + \mu_2d_2 + \mu_3d_3.
     \end{equation}
     The right-hand side is clearly nonnegative. If $(a_{00}-\mu_0) = 0$, then $\mu_1 =\mu_2=\mu_3 = 0$.  If $(a_{00}-\mu_0) > 0$, then the minimality of $a_{00}$ implies $\mu_0 = 0$.
\end{proof}

\begin{lemma}\label{lemma2cases}
     If there is no $m \in \{0,1,2,3\}$ such that $a_{ii} = b_{m,ii}$ for at least two distinct indices $i \in \{0,1,2,3\} \setminus \{m\}$, then there exists an arrangement $\{i,j,k,l\} = \{0,1,2,3\}$ such that $a_{ii}d_i = a_{jj}d_j \neq a_{kk}d_k = a_{ll}d_l$, and $\varphi^{-1}(a_{ii}d_i) \cup \varphi^{-1}(a_{kk}d_k) =  \{(a_{00},0,0,0), (0,a_{11},0,0),(0,0,a_{22},0),(0,0,0,a_{33})\}$.
\end{lemma}

\begin{proof}
    Suppose that for every $m \in \{0,1,2,3\}$ we have $a_{ii} < b_{m,ii}$ for at least two distinct indices $i \in \{0,1,2,3\} \setminus \{m\}$ (note that $a_{ii} \leq b_{m,ii}$ always holds by definition). 
    
    Suppose that there exists $z \in \varphi^{-1}(a_{00} d_0)$ with $|\textup{supp}(z)| \geq 2$. By Lemma \ref{lemma_catenary_minrelations}, $0 \notin \textup{supp}(z)$. Without loss of generality, suppose that $z=(0,a_{01},a_{02},a_{03})$ where $a_{01},a_{02}>0$. Then, $a_{00} = b_{1,00} = b_{2,00}$. This implies $a_{11} < b_{2,11}$, $a_{22} < b_{1,22}$, $a_{33} < b_{1,33}$, and $a_{33} <b_{2,33}$. The last two inequalities give $\varphi^{-1}(a_{33}d_3) = \{(0,0,0,a_{33}),(a_{30},0,0,0)\}$, so $a_{30}>0$ and hence $a_{33} = b_{0,33}$. This implies $a_{11} < b_{0,11}$ and $a_{22} < b_{0,22}$. Now, since $a_{11} < b_{0,11}$ and $a_{11} < b_{2,11}$, we have $\varphi^{-1}(a_{11}d_1) = \{(0,a_{11},0,0),(0,0,0,a_{13})\}$, so $a_{13} > 0$ and hence $a_{11} = b_{3,11}$. Similarly, $\varphi^{-1}(a_{22}d_2) = \{(0,0,a_{22},0),(0,0,0,a_{23})\}$ and $a_{22} = b_{3,22}$, which is a contradiction, as now there exists at most one $i \in \{0,1,2\}$ such that $a_{ii} < b_{3,ii}$. 

    Hence, for every $z \in \varphi^{-1}(a_{00} d_0)$, we have $|\textup{supp}(z)| = 1$.
    Proceeding analogously with $a_{11}d_1$, $a_{22}d_2$, and $a_{33}d_3$, we obtain that for every $i \in \{0,1,2,3\}$ and every $z \in \varphi^{-1}(a_{ii} d_i)$, we have $|\textup{supp}(z)| = 1$.

    Without loss of generality, suppose that $a_{00}d_0 = a_{01}d_1$.
    Observe that $a_{01} \geq a_{11}$. If $a_{01} > a_{11}$, then
    \begin{equation}\label{notprop-eq1}
        a_{00}d_0 = a_{01}d_1 = ( a_{01}-a_{11})d_1 + a_{1i}d_i
    \end{equation}
    for some $i \in \{0,2,3\}$. Clearly $a_{1i} > 0$. If $i=0$, subtracting $d_0$ from both sides of \eqref{notprop-eq1} contradicts the minimality of $a_{00}$. If $i \neq 0$, we get a factorization $z \in  \varphi^{-1}(a_{00} d_0)$ with $|\textup{supp}(z)| = 2$, which is a contradiction. 
    
    Hence $a_{01} = a_{11}$, and therefore $a_{00}d_0 = a_{11}d_1$. Analogously, we get $a_{22}d_2 = a_{ii}d_i$ for some $i \in \{0,1,3\}$. If $i \in \{0,1\}$, then $a_{00}d_0 = a_{11}d_1 = a_{22}d_2$, so $a_{11} = b_{0,11}$ and $a_{22} = b_{0,22}$, which is contrary to our assumption. Hence $a_{22}d_2 = a_{33}d_3 \neq a_{00}d_0 = a_{11}d_1$, and the claim follows.
\end{proof}

If there exists $m \in \{0,1,2,3\}$ such that $a_{ii} = b_{m,ii}$ for at least two distinct indices $i \in \{0,1,2,3\} \setminus \{m\}$, we call $\langle d_0,d_1,d_2,d_3 \rangle$ \emph{primary}. Moreover, for a primary semigroup $\langle d_0,d_1,d_2,d_3 \rangle$, unless stated otherwise, we always assume throughout this paper that the generators are ordered so that $a_{11} = b_{0,11}$ and $a_{22} = b_{0,22}$. 

Numerical semigroups with embedding dimension four that are not primary will be referred to as \emph{secondary} semigroups. By Lemma \ref{lemma2cases}, if $\langle d_0,d_1,d_2,d_3 \rangle$ is a secondary semigroup, then there exists an arrangement $\{i,j,k,l\} = \{0,1,2,3\}$ such that $a_{ii}d_i = a_{jj}d_j \neq a_{kk}d_k = a_{ll}d_l$, and $\varphi^{-1}(a_{ii}d_i) \cup \varphi^{-1}(a_{kk}d_k) =  \{(a_{00},0,0,0), (0,a_{11},0,0), (0,0,a_{22},0),(0,0,0,a_{33})\}$. Moreover, for a secondary semigroup $\langle d_0,d_1,d_2,d_3 \rangle$, unless stated otherwise, we always assume throughout this paper that the generators are ordered so that $\varphi^{-1}(a_{00}d_0) = \varphi^{-1}(a_{11}d_1)= \{(a_{00},0,0,0), (0,a_{11},0,0)\}$ and $ \varphi^{-1}(a_{22}d_2) = \varphi^{-1}(a_{33}d_3) =  \{(0,0,a_{22},0), (0,0,0,a_{33})\}$. 

A numerical semigroup with embedding dimension four, $S=\langle d_0,d_1,d_2,d_3 \rangle$, is said to be \emph{ordered} if its generators are arranged according to the ordering defined above --- that is, the ordering assumed for primary semigroups if $S$ is primary, and the ordering assumed for secondary semigroups if $S$ is secondary. We adopt this convention so that certain definitions and constructions can be stated uniformly for any ordered semigroup $\langle d_0,d_1,d_2,d_3 \rangle$, without requiring a prior assumption on whether it is primary or secondary.

\subsection{Geometric framework}\label{subsection_framework}

From this point on, we will focus on the Ap{\'e}ry set of $\langle d_0,d_1,d_2,d_3 \rangle$ with respect to $d_0$, which we henceforth refer to simply as the Ap{\'e}ry set. For brevity, we also write $b_{0,ij} = b_{ij}$ for all $i,j \in \{0,1,2,3\}$. 


Consider the unit cubes in $\mathbb{R}^3$. Each unit cube $[i,i+1] \times [j,j+1] \times [k,k+1] \subset \mathbb{R}^3$ has natural coordinates $(i,j,k) \in \mathbb{N}^3$ and is labeled with the value $id_1 + jd_2 + kd_3$. We will refer to the collection of all such labeled cubes as \emph{the initial collection of cubes}. We denote the unit cube $[i,i+1] \times [j,j+1] \times [k,k+1]$ by $[[i,j,k]]$.
For a point $P =(a,b,c) \in \mathbb{Z}^3$, we call the region $\{x > a, y>b, z>c\}$ \emph{the region associated to} $P$. We write $[[a,b,c]] \geq [[a',b',c']]$ when $a \geq a'$, $b \geq b'$, $c \geq c'$, and $[[a,b,c]] > [[a',b',c']]$ when $[[a,b,c]] \geq [[a',b',c']]$ and $[[a,b,c]] \neq [[a',b',c']]$.


\begin{lemma}\label{lemma_procedure_1}
   No cube labeled with an element of the Ap{\'e}ry set is in the region associated to a point $(\mu_1,\mu_2,\mu_3) \in \mathbb{Z}^3$ satisfying $\mu_1d_1 + \mu_2d_2 + \mu_3d_3 = \mu d_0$ for some positive integer $\mu$. 
\end{lemma}

\begin{proof}
    Recall that the Ap{\'e}ry set with respect to $d_0$ is the set of the smallest nonnegative integers from each residue class modulo $d_0$ that are in the numerical semigroup $\langle d_1,d_2,d_3\rangle$. Suppose to the contrary, that for some residue $r \in \{0,1,\dots,d_0-1\}$ there is a cube $[[\lambda_1,\lambda_2,\lambda_3]]$ labeled with the smallest number $\lambda d_0 + r$ ($\lambda \in \mathbb{N}$) that is in $\langle d_1,d_2,d_3\rangle$ such that $[[\lambda_1,\lambda_2,\lambda_3]] \geq [[\mu_1,\mu_2,\mu_3]]$. This gives
    \begin{equation}\label{lemma_procedure_1_eq}
     (\lambda_1 -\mu_1) d_1 + (\lambda_2-\mu_2) d_2 + (\lambda_3-\mu_3) d_3 = (\lambda-\mu) d_0 + r
    \end{equation}
    where all coefficients on the left-hand side are nonnegative. This means that the value in \eqref{lemma_procedure_1_eq} is congruent to $r$ modulo $d_0$, is nonnegative, is in $\langle d_1,d_2,d_3\rangle$, and is strictly less than $\lambda d_0 + r$, which contradicts the minimality of $\lambda d_0 + r$.
\end{proof}


\begin{thm}\label{theoremprocedure}
    Delete from the initial collection of cubes the regions $\{x > b_{11}\}$, $\{y > b_{22}\}$, and $\{z > b_{33}\}$.
    Next, for every point $(a_{01},a_{02},a_{03}) \in \mathbb{N}^3$ satisfying
     \begin{equation}\label{minrelation-0}
          a_{01} d_1 + a_{02}d_2 + a_{03} d_3 = a_{00} d_0
     \end{equation}
    delete the region associated to it. 
    Furthermore, for all identities of the form
     \begin{equation}\label{theoremprocedure_eq1}
         \lambda_0d_0 + \lambda_1d_1=\lambda_2d_2 + \lambda_3d_3
     \end{equation}
    with $\lambda_i \in \mathbb{Z}_{>0}$, $\lambda_2 < b_{22}$, and $\lambda_3 < b_{33}$, delete the region associated to $(-\lambda_1,\lambda_2,\lambda_3)$. For all identities of the form
    \begin{equation}\label{theoremprocedure_eq2}
         \lambda_0d_0 + \lambda_2d_2=\lambda_1d_1 + \lambda_3d_3
    \end{equation}
    with $\lambda_i \in \mathbb{Z}_{>0}$, $\lambda_1 < b_{11}$, and $\lambda_3 < b_{33}$, delete the region associated to $(\lambda_1,-\lambda_2,\lambda_3)$. Finally, for all identities of the form
    \begin{equation}\label{theoremprocedure_eq3}
        \lambda_0d_0 + \lambda_3d_3=\lambda_1d_1 + \lambda_2d_2
    \end{equation}
    with $\lambda_i \in \mathbb{Z}_{>0}$, $\lambda_1 < b_{11}$, and $\lambda_2 < b_{22}$, delete the region associated to $(\lambda_1,\lambda_2,-\lambda_3)$.
    Then, the remaining collection of cubes is the set of all cubes from the initial collection of cubes that are labeled with an element of the Ap{\'e}ry set.
\end{thm}

\begin{proof}
Let $R$ be the remaining collection of cubes.
The condition $\lambda_i < b_{ii}$ ensures that only the relevant identities are considered, since we delete the regions $\{x > b_{11}\}$, $\{y > b_{22}\}$, and $\{z > b_{33}\}$. Moreover, by Lemma \ref{lemma_procedure_1}, every cube labeled with an element of the Ap{\'e}ry set is in $R$, since deleting the regions $\{x > b_{11}\}$, $\{y > b_{22}\}$, and $\{z > b_{33}\}$ is the same as deleting the regions associated to $(b_{11}, - b_{12},-b_{13})$, $(-b_{21},b_{22},-b_{23})$, and $(-b_{31},-b_{32},b_{33})$. 

Therefore we only need to prove that $R$ contains only cubes labeled with an element of the Ap{\'e}ry set. Suppose to the contrary, that for some remainder $r \in \{0,1,\dots,d_0-1\}$, $\lambda d_0 + r$ and $\lambda' d_0 + r$, where $\lambda' > \lambda$, are labels in $R$. This means that
    \begin{equation}\label{eq1theorem}
        \lambda d_0 + r = \lambda_1d_1 + \lambda_2d_2 + \lambda_3d_3
    \end{equation}
     \begin{equation}\label{eq2theorem}
        \lambda' d_0 + r = \lambda_1'd_1 + \lambda_2'd_2 + \lambda_3'd_3
    \end{equation}
    where $[[\lambda_1,\lambda_2,\lambda_3]], [[\lambda_1',\lambda_2',\lambda_3']] \in R$.
    Subtracting \eqref{eq1theorem} from \eqref{eq2theorem} we get
    \begin{equation}\label{eq3theorem}
        (\lambda' - \lambda)d_0 = (\lambda_1' - \lambda_1)d_1 + (\lambda_2' - \lambda_2)d_2 + (\lambda_3' - \lambda_3)d_3
    \end{equation}
    where $(\lambda' - \lambda)>0$. Now, we consider cases of how many coefficients on the right-hand side of \eqref{eq3theorem} are positive. The left-hand side is positive; hence, at least one of the coefficients on the right-hand side is positive. 
    
If exactly one coefficient is positive, without loss of generality, suppose that $(\lambda_1' - \lambda_1)$ is the one. Moving the non-positive integers to the other side gives
    \begin{equation}\label{eq4theorem}
        (\lambda' - \lambda)d_0  + (\lambda_2 - \lambda_2')d_2 + (\lambda_3 - \lambda_3')d_3= (\lambda_1' - \lambda_1)d_1.
    \end{equation} 
Since all the above coefficients are nonnegative and $(\lambda_1' - \lambda_1)$ and $(\lambda' - \lambda) $ are positive, the minimality of $b_{11}$ gives $b_{11} \leq (\lambda_1' - \lambda_1) \leq \lambda_1'$, which contradicts $[[\lambda_1',\lambda_2',\lambda_3']] \in R$.

If exactly two coefficients on the right-hand side of \eqref{eq3theorem} are positive, without loss of generality, suppose that $(\lambda_1' - \lambda_1)$ and $(\lambda_2' - \lambda_2)$ are the two. Moving $(\lambda_3' - \lambda_3)$ to the other side gives
    \begin{equation}\label{eq5theorem}
        (\lambda' - \lambda)d_0  + (\lambda_3 - \lambda_3')d_3= (\lambda_1' - \lambda_1)d_1  + (\lambda_2' - \lambda_2)d_2.
    \end{equation}
The above identity satisfies the conditions of the current theorem, thus the region $\{x > (\lambda_1' - \lambda_1), y > (\lambda_2' - \lambda_2)\}$ was deleted. Again, this contradicts $[[\lambda_1',\lambda_2',\lambda_3']] \in R$, since $\lambda'_1 \geq (\lambda_1' - \lambda_1)$ and $\lambda'_2 \geq (\lambda_2' - \lambda_2)$.

Now, the only case left is where all the coefficients in \eqref{eq3theorem} are positive. We substitute $\mu =(\lambda' - \lambda)  > 0$ and $\mu_i =(\lambda_i' - \lambda_i)  >0$ for $i = 1,2,3$. We have
    \begin{equation}\label{eq6theorem}
         \mu d_0 = \mu_1d_1 + \mu_2d_2 + \mu_3d_3
    \end{equation}
and since $[[\mu_1,\mu_2,\mu_3]] \leq [[\lambda_1',\lambda_2',\lambda_3']] \in R$, we have $[[\mu_1,\mu_2,\mu_3]] \in R$. 
From the minimality of $a_{00}$, we get $a_{00} \leq \mu$. Subtracting $a_{00}d_0 = a_{01}d_1 + a_{02}d_2 + a_{03}d_3$ from \eqref{eq6theorem} gives
\begin{equation}\label{eq7theorem}
    (\mu - a_{00}) d_0 = (\mu_1 - a_{01})d_1 +  (\mu_2 - a_{02})d_2 + (\mu_3 - a_{03})d_3.
\end{equation}
    
If $\mu = a_{00}$, then $(0,\mu_1,\mu_2,\mu_3) \in \varphi^{-1}(a_{00}d_0)$, hence the region associated to $(\mu_1,\mu_2,\mu_3)$ had been deleted, which contradicts $[[\mu_1,\mu_2,\mu_3]] \in R$. 

Hence $(\mu - a_{00}) > 0$. Now, we repeat the same case by case consideration with \eqref{eq7theorem} as we did with \eqref{eq3theorem}, which results in all the coefficients in \eqref{eq7theorem} being positive. However, this implies that the cube $[[\mu_1,\mu_2,\mu_3]]$ is in the region associated to $(a_{01}, a_{02}, a_{03})$, which contradicts $[[\mu_1,\mu_2,\mu_3]] \in R$. 
\end{proof}

Throughout this paper, $R$ denotes the collection of all cubes from the initial collection of cubes that are labeled with an element of the Ap{\'e}ry set.
We will refer to equations of the forms \eqref{theoremprocedure_eq1}, \eqref{theoremprocedure_eq2}, \eqref{theoremprocedure_eq3} (where $\lambda_i \in \mathbb{Z}_{>0}$) as the \emph{$(2,2)$-type relations}. In addition, we will refer to points that are constructed from the $(2,2)$-type relations (as in Theorem \ref{theoremprocedure}) as the \emph{$(2,2)$-type points}. 

\subsection{$L$-shapes}\label{subsection_Lshapes}

Define an \emph{$L$-shape} to be a subset of the initial collection of cubes such that:
\begin{itemize}
    \item all its labels are elements of the Ap{\'e}ry set,
    \item every element of the Ap{\'e}ry set labels exactly one cube in it,
    \item if a cube $[[a,b,c]]$, $(a,b,c) \in \mathbb{N}^3$, is not in it, then neither is any cube in the region associated to $(a,b,c)$.
\end{itemize}
This notion is adopted from \cite{AGUILOGOST2015} and is established in the literature.

\begin{lemma}\label{lemma_procedure_2}
     Delete from $R$ one of the regions $\{x > a_{11}\}$, $\{y > a_{22}\}$, or $\{z > a_{33}\}$. Then every element of the Ap{\'e}ry set still labels some remaining cube.
\end{lemma}

\begin{proof}
Delete the region $\{x > a_{11}\}$ from $R$ (the other two cases follow analogously). 
If $a_{11} = b_{11}$, there is nothing to prove; thus, suppose $a_{11} < b_{11}$. Suppose to the contrary, that there is some element of the Ap{\'e}ry set that does not label any cube in $R \cap \{x \leq a_{11}\}$. By Theorem \ref{theoremprocedure}, then there exists $[[\lambda_1, \lambda_2,\lambda_3]] \in R$ with $\lambda_1 \geq a_{11}$ that is labeled with the element in question. There exists a positive integer $\lambda$ such that $a_{11} > (\lambda_1 - \lambda a_{11}) \geq 0$. Using $a_{11}d_1 = a_{12}d_2 + a_{13}d_3$, where $a_{12},a_{13} \in \mathbb{N}$ (recall $a_{11} < b_{11}$), we obtain
    \begin{equation}
        \lambda_1d_1 + \lambda_2d_2 + \lambda_3d_3 = (\lambda_1 - \lambda a_{11})d_1 + (\lambda_2+\lambda a_{12})d_2 + (\lambda_3+\lambda a_{13})d_3.
    \end{equation}
The cube $[[\lambda_1 - \lambda a_{11},\lambda_2+\lambda a_{12},\lambda_3+\lambda a_{13} ]]$ has the same label as $[[\lambda_1, \lambda_2,\lambda_3]]$, hence it is in $R$ by Theorem \ref{theoremprocedure}. Since $a_{11} >(\lambda_1 - \lambda a_{11})$, it is in $R \cap \{x \leq a_{11}\}$ --- contradiction.
\end{proof}

\begin{lemma}\label{lemma_procedure_3}
    Delete from $R$ the regions $\{x > a_{11}\}$, $\{y > a_{22}\}$, and $\{z > a_{33}\}$. Then, in the resulting collection, each element of the Ap{\'e}ry set labels at most one cube.
\end{lemma}
\begin{proof}
     Suppose to the contrary, that for two distinct $[[\lambda_1, \lambda_2, \lambda_3]], [[\lambda_1', \lambda_2', \lambda_3']] \in R \cap \{x \leq a_{11}, y \leq a_{22}, z \leq a_{33} \}$, we have
    \begin{equation*}
        \lambda_1d_1 + \lambda_2d_2 + \lambda_3d_3 = \lambda_1'd_1 + \lambda_2'd_2 + \lambda_3'd_3 
    \end{equation*}
    \begin{equation*}
        \implies  (\lambda_1' - \lambda_1)d_1 + (\lambda_2' - \lambda_2)d_2 + (\lambda_3' - \lambda_3)d_3 = 0.
    \end{equation*}
    If all coefficients above are zero, we have nothing to prove. Therefore either one or two coefficients above are positive. Regardless, without loss of generality, we can suppose that
    \begin{equation}\label{propeq3}
         (\lambda_1' - \lambda_1)d_1 = (\lambda_2 - \lambda_2')d_2 + (\lambda_3 - \lambda_3')d_3
    \end{equation}
    where every coefficient above is nonnegative. However, $(\lambda_1' - \lambda_1)$ is positive, because otherwise all coefficients are zero. This implies $a_{11} \leq (\lambda_1' - \lambda_1) \leq \lambda_1'$, which contradicts $[[\lambda_1', \lambda_2', \lambda_3']] \in R \cap \{x \leq a_{11}, y \leq a_{22}, z \leq a_{33} \}$.
\end{proof}

\begin{prop}\label{prop_primary}
    If $\langle d_0,d_1,d_2,d_3 \rangle$ is primary, then $R \cap \{z \leq a_{33}\}$ is an $L$-shape.
\end{prop}
\begin{proof}
    If $a_{ii} = b_{ii}$ for $i = 1,2,3$, then $R = R \cap \{x \leq a_{11}, y \leq a_{22}, z \leq a_{33} \}$, hence Lemma \ref{lemma_procedure_3} yields the claim. Otherwise, we have $a_{ii} < b_{ii}$ for exactly one $i \in \{1,2,3\}$. In such case, Lemmas \ref{lemma_procedure_2} and \ref{lemma_procedure_3} yield the claim.
\end{proof}


\begin{prop}\label{prop_secondary}
    Suppose that for $\langle d_0,d_1,d_2,d_3 \rangle$ we have $a_{11}=b_{11}$ and $a_{22}d_2 = a_{33}d_3$. Then there is a bijection $f : R \cap \{y \geq a_{22}\} \to R \cap \{z \geq a_{33}\}$ sending $[[\lambda_1,\lambda_2,\lambda_3]] \mapsto [[\lambda_1,\lambda_2-a_{22},\lambda_3+a_{33}]]$. Moreover, $R \cap \{y \leq a_{22}\}$ and $R \cap \{z \leq a_{33}\}$ are $L$-shapes.
\end{prop}

\begin{proof}
    For $[[\lambda_1,\lambda_2,\lambda_3]] \in R \cap \{y \geq a_{22}\} $, the cube $[[\mu_1,\mu_2 - a_{22},\mu_3 + a_{33}]]$ is in $R$ by Theorem \ref{theoremprocedure}, because it has the same label as $[[\lambda_1,\lambda_2,\lambda_3]]$, so it is labeled with an element of the Ap{\'e}ry set. Thus, the function $f: [[\lambda_1,\lambda_2,\lambda_3]] \mapsto [[\lambda_1,\lambda_2-a_{22},\lambda_3+a_{33}]]$ maps $R \cap \{y \geq a_{22}\}$ to $R \cap \{z \geq a_{33}\}$ and is also clearly injective. Similarly, we obtain the injection $f^{-1} : R \cap \{z \geq a_{33}\} \to R \cap \{y \geq a_{22}\} $ sending $[[\lambda_1,\lambda_2,\lambda_3]] \mapsto [[\lambda_1,\lambda_2+a_{22},\lambda_3-a_{33}]]$, which yields that $f$ is a bijection.   
    
    Let $T = R \cap \{z \leq a_{33}\}$. Since $f$ preserves labels, every element of the Ap{\'e}ry set labels a cube in $T$. Thus, to prove that $T$ is an $L$-shape we only have to prove that it has distinct labels.  
     
    Otherwise, we have
    \begin{equation}\label{propnottrick-eq1}
        \mu_1d_1 + \mu_2d_2 + \mu_3d_3 = \nu_1d_1 + \nu_2d_2 + \nu_3d_3
    \end{equation}
    where $[[\mu_1,\mu_2,\mu_3]]$ and $[[\nu_1,\nu_2,\nu_3]]$ are two distinct cubes in $T$. Without loss of generality, suppose $\mu_2 \geq \nu_2$. 
    
    If $\mu_1 \leq \nu_1$ and $\mu_3 \leq \nu_3$, then
    \begin{equation}\label{propnottrick-eq2}
         (\mu_2-\nu_2)d_2  = (\nu_1-\mu_1)d_1  + (\nu_3-\mu_3)d_3
    \end{equation}
    where all coefficients above are nonnegative, which implies $(\mu_2-\nu_2) \geq a_{22}$ (as $[[\mu_1,\mu_2,\mu_3]] \neq [[\nu_1,\nu_2,\nu_3]]$). If $\nu_3 \geq (\nu_3-\mu_3) \geq a_{33}$, we get a contradiction with $[[\nu_1,\nu_2,\nu_3]] \in T$, hence $(\nu_3-\mu_3) < a_{33}$. This gives
    \begin{equation}\label{propnottrick-eq21}
         (\mu_2-\nu_2-a_{22})d_2 + (a_{33}-\nu_3+\mu_3)d_3  = (\nu_1-\mu_1)d_1
    \end{equation}
    where all coefficients above are nonnegative and $(a_{33}-\nu_3+\mu_3)$ is positive. Hence $(\nu_1-\mu_1)$ is also positive, which implies $\nu_1 \geq (\nu_1-\mu_1) \geq a_{11} = b_{11}$, contradicting $[[\nu_1,\nu_2,\nu_3]] \in T$.

    If $\mu_1 > \nu_1$ and $\mu_3 \leq \nu_3$, then
    \begin{equation}\label{propnottrick-eq3}
         (\mu_2-\nu_2)d_2 + (\mu_1-\nu_1)d_1  =  (\nu_3-\mu_3)d_3
    \end{equation}
    where all coefficients above are nonnegative and $(\mu_1-\nu_1)$ is positive, which implies $\nu_3 \geq (\nu_3-\mu_3) \geq a_{33}$, contradicting $[[\nu_1,\nu_2,\nu_3]] \in T$.

    Lastly, if $\mu_1 \leq \nu_1$ and $\mu_3 > \nu_3$, then
    \begin{equation}\label{propnottrick-eq4}
         (\mu_2-\nu_2)d_2 + (\mu_3-\nu_3)d_3  =  (\nu_1-\mu_1)d_1
    \end{equation}
     where all coefficients above are nonnegative and $(\mu_3-\nu_3)$ is positive, which implies $\nu_1 \geq (\nu_1-\mu_1) \geq a_{11} = b_{11}$, so again a contradiction with $[[\nu_1,\nu_2,\nu_3]] \in T$.

    Analogously, we obtain that $R \cap \{y \leq a_{22}\}$ is an $L$-shape.
\end{proof}



In particular, since secondary semigroups satisfy $a_{11}=b_{11}$ and $a_{22}d_2 = a_{33}d_3$ by definition, Proposition \ref{prop_secondary} applies to all secondary semigroups.




\begin{ex}\label{example_T}
Consider $\langle d_0,d_1,d_2,d_3 \rangle = \langle 196,225,256,289 \rangle$, which has embedding dimension four. One can check that
\begin{alignat*}{2}
    &a_{00}d_0 = 17 \cdot 196 &&= 2 \cdot 225 + 9 \cdot 256 + 2 \cdot 289,\\
    &a_{11}d_1 = 7 \cdot 225 &&= 1 \cdot 196 + 2 \cdot 256 +3 \cdot 289,\\
   &a_{22}d_2 =15 \cdot 256 &&= 15 \cdot 196 + 4 \cdot 225,\\
   &a_{33}d_3 =4 \cdot 289 &&= 4 \cdot 225 + 1 \cdot 256.
\end{alignat*}
Hence $\langle 196,225,256,289 \rangle$ is primary. Figure \ref{fig_theorem_primary} shows the corresponding collection $R$. The cubes colored red are in the region $R \cap \{z \geq a_{33}\}$ and the cubes colored gray form the $L$-shape from Proposition \ref{prop_primary}. 
Next, consider $\langle d_0,d_1,d_2,d_3 \rangle =\langle 145,203,74,111 \rangle$, which has embedding dimension four. One can check that
\begin{alignat*}{2}
    & \varphi^{-1}(a_{00}d_0) = \varphi^{-1}(7 \cdot 145) &&= \{(7,0,0,0),(0,5,0,0)\},\\
    &\varphi^{-1}(a_{11}d_1) =\varphi^{-1}(5 \cdot 203) &&= \{(7,0,0,0),(0,5,0,0)\},\\
   &\varphi^{-1}(a_{22}d_2) = \varphi^{-1}( 3\cdot 74) &&= \{(0,0,3,0),(0,0,0,2)\},\\
   &\varphi^{-1}(a_{33}d_3) = \varphi^{-1}(2 \cdot 111)  &&= \{(0,0,3,0),(0,0,0,2)\}.
\end{alignat*}
Hence $\langle 145,203,74,111 \rangle$ is secondary.  Figure \ref{fig_theorem_secondary} shows the corresponding collection $R$. Similarly, red marks the region $R \cap \{z \geq a_{33}\}$, while gray marks an $L$-shape from Proposition \ref{prop_secondary}. 
\end{ex}


\begin{figure}[h]
    \centering
    \begin{minipage}[t]{0.48\textwidth}
        \centering
        \includegraphics[width=0.7\textwidth]{tintar_theorem_primary1.pdf}
        \caption{}
        \label{fig_theorem_primary}
    \end{minipage}
    \begin{minipage}[t]{0.48\textwidth}
        \centering
        \includegraphics[width=0.7\textwidth]{tintar_theorem_secondary.pdf}
        \caption{}
        \label{fig_theorem_secondary}
    \end{minipage}
\end{figure}


\section{Minimal presentations}\label{section_minimal_presentations}

In this section, we investigate the set of Betti elements for primary and secondary semigroups. We start with preliminary definitions and then proceed to the main results. Theorems \ref{theorembettifind_primary} and \ref{theorembettifind_secondary} relate Betti elements to an $L$-shape for primary and secondary semigroups respectively. Lemmas \ref{lemma_catenary_primary} and \ref{lemma_catenarycase1_primary} determine the sets of factorizations of Betti elements arising from $(2,2)$-type points in the primary case, while Lemmas \ref{lemma_catenary_lambda1_secondary} and \ref{lemma_catenary_lambda23_secondary} establish the analogous results for secondary semigroups. Finally, Lemma \ref{lemma_catenaryT_primary} characterizes the $(2,2)$-type points which do not correspond to Betti elements in the primary case, while Lemmas \ref{lemma_catenaryT_secondary_case1} and \ref{lemma_catenaryT_secondary_case2} do the same in the secondary case. This framework will be used in Section \ref{section_type} to study minimal presentations from a geometric perspective.

Bresinsky’s work \cite{Bresinsky1988}, which underlies the bound in Theorem \ref{theoremtypebre}, contains results closely related to those of this section (cf.\ \cite[Theorem 3]{Bresinsky1988}). Our approach differs in two main respects. First, in the classification of numerical semigroups with $e(S)=4$, we distinguish primary and secondary semigroups, whereas Bresinsky introduced three families. Second, we relate Betti elements and minimal presentations to the Ap{\'e}ry set via an $L$-shape, allowing us to make use of its geometric structure, whereas Bresinsky's approach was commutative algebra centered.



Let $S$ be minimally generated by $\{d_0,d_1,\dots,d_k\}$.
Two elements $z$ and $z'$ of $\mathbb{N}^{k+1}$ are said to be \emph{$\mathcal{R}$-related} if there exists a finite sequence $z = z_1, \dots, z_p = z'$ in $\mathbb{N}^{k+1}$ such that $\text{supp}(z_i) \cap \text{supp}(z_{i+1})$ is nonempty for all $i \in \{1, \dots, p-1\}$. We write $z \, \mathcal{R} \, z'$ in this case. We refer to the equivalence classes of $\varphi^{-1}(s)$ under the relation $\mathcal{R}$ as the \emph{$\mathcal{R}$-classes}. 
An element $s \in S$ is called a \emph{Betti element} if there exist two elements of $\varphi^{-1}(s)$ that are not $\mathcal{R}$-related (i.e.\ $\varphi^{-1}(s)$ has at least two $\mathcal{R}$-classes). The set of all Betti elements of $S$ is denoted $\text{Betti}(S)$.






Let $\mathcal{R}_1, \dots, \mathcal{R}_l$ be all the $\mathcal{R}$-classes of $\varphi^{-1}(b)$ for some $b \in \text{Betti}(S)$. We choose $v_i \in \mathcal{R}_i$ for each $i \in \{1,\dots,l\}$. Let $\rho_b$ be a set of $l-1$ pairs of elements from $\{v_1, \dots, v_l\}$ such that the graph with vertices $\{v_1, \dots, v_l\}$, having an edge between $v_i$ and $v_j$ if and only if $(v_i,v_j)$ or $(v_j,v_i)$ is in $\rho_b$, is connected. 
\begin{thm}\label{theorem_minimalpresentation}
    \textup{\cite{Rosales1996MinimalRelation}}
    The set $\rho = \bigcup_{b \in \text{Betti}(S)} \rho_b$ is a minimal presentation of $S$ and every minimal presentation of $S$ is of this form.
\end{thm}





For the rest of this section, let $S = \langle d_0,d_1,d_2,d_3 \rangle$ be an ordered numerical semigroup with embedding dimension four. Let $T = R \cap \{z \leq a_{33}\}$. By Propositions \ref{prop_primary} and \ref{prop_secondary}, $T$ is an $L$-shape.
Let $V \subset \mathbb{Z}^3$ be a minimal set (with respect to inclusion) of points $(x,y,z)$ such that $xd_1 + yd_2 +zd_3$ is a nonnegative multiple of $d_0$, and such that removing from the initial collection of cubes the regions associated to the points of $V$ yields $T$. 
Define
$$ U = \{ \max\{ x,0\} d_1 + \max\{y,0\}d_2 + \max\{z,0\}d_3 \mid (x,y,z) \in V \}.$$
Note that $U$ is independent of the choice of $V$. Furthermore, assume that every point $(x,y,z) \in V$ such that
$$ \max\{ x,0\} d_1 + \max\{y,0\}d_2 + \max\{z,0\}d_3 \notin \{(a_{00}d_0,b_{11}d_1,b_{22}d_2,a_{33}d_3)\}$$
is a $(2,2)$-type point. Such $V$ exists by Theorem \ref{theoremprocedure}. Throughout this paper, we will always assume that $V$ satisfies this.

\begin{remark}
   The above assumption may fail for certain choices of $V$. For example, if $a_{00}d_0 = a_{11}d_1$ and there exists $(-\lambda_1,\lambda_2,\lambda_3) \in V$ such that $\lambda_1 = \mu a_{11}$ for some $\mu \in \mathbb{Z}_{>0}$, then we could take $(0,\lambda_2,\lambda_3)$ as an element of $V$ instead of $(-\lambda_1,\lambda_2,\lambda_3)$, since $\lambda_2d_2 + \lambda_3d_3 = (\lambda_0 + \mu a_{00})d_0$.
\end{remark}

For the semigroups $\langle 196,225,256,289 \rangle$ and $\langle 145,203,74,111 \rangle$ from Example \ref{example_T}, shown respectively in Figures \ref{fig_betti_primary} and \ref{fig_betti_secondary}, the cubes shown in gray form the corresponding collections $T$. The corresponding sets $U$ are the set of labels of the cubes shown in red.

\begin{figure}[h]
    \centering
    \begin{minipage}[t]{0.48\textwidth}
        \centering
        \includegraphics[width=0.7\textwidth]{tintar_betti_primary.pdf}
        \caption{}
        \label{fig_betti_primary}
    \end{minipage}
    \begin{minipage}[t]{0.48\textwidth}
        \centering
        \includegraphics[width=0.7\textwidth]{tintar_betti_secondary.pdf}
        \caption{}
        \label{fig_betti_secondary}
    \end{minipage}
\end{figure}



\subsection{Primary case}

\begin{lemma}\label{lemmabettifind_primary}
  Suppose that $S$ is primary. If $b \in \textup{Betti}(S)$ labels a cube in $R$, then $b \in U$.
\end{lemma}
    
\begin{proof}
    In this proof, it will be more convenient to drop the assumption that $S$ is ordered --- we only suppose that $a_{ii}=b_{ii}$ for at least two distinct $i \in \{1,2,3\}$ and that  $T = R \cap \{x \leq a_{11},y\leq a_{22}, z \leq a_{33}\}$.
    
    Suppose that $b$ labels a cube in $R$. Then $b \in \textup{Ap}(S,d_0)$ by Theorem \ref{theoremprocedure}. Since $T$ is an $L$-shape, $b$ labels a cube in $T$.   Consider the factorization $z_0 =(0,\lambda_1,\lambda_2, \lambda_3) \in \varphi^{-1}(b)$ with $[[\lambda_1,\lambda_2,\lambda_3]] \in T$. Let $(\lambda_0',\lambda_1',\lambda_2', \lambda_3')$ be a different factorization of $b$. We have
    \begin{equation}\label{betticharacter_eq0}
        (\lambda_1-\lambda_1')d_1 +  (\lambda_2-\lambda_2')d_2 +  (\lambda_3-\lambda_3')d_3 = \lambda_0'd_0.
    \end{equation}

    If all coefficients on the left-hand side are nonnegative, then $\lambda_0'>0$, as the two factorizations are distinct. This contradicts $[[\lambda_1,\lambda_2,\lambda_3]] \in T$ by Lemma \ref{lemma_procedure_1}. 


    Suppose that exactly one coefficient on the left-hand side of \eqref{betticharacter_eq0} is nonnegative. Without loss of generality, suppose that every coefficient below is nonnegative.
    \begin{equation}\label{betticharacter_eq00}
        (\lambda_1-\lambda_1')d_1   = \lambda_0'd_0 + (\lambda_2'-\lambda_2)d_2 +  (\lambda_3'-\lambda_3)d_3.
    \end{equation}
    This gives $\lambda_1 \geq (\lambda_1-\lambda_1') \geq a_{11}$, which contradicts $[[\lambda_1,\lambda_2,\lambda_3]] \in T$.
    

    Suppose that exactly two coefficients on the left-hand side of \eqref{betticharacter_eq0} are nonnegative. Without loss of generality, suppose that every coefficient below is nonnegative.
     \begin{equation}\label{betticharacter_eq01}
        (\lambda_1-\lambda_1')d_1 +  (\lambda_2-\lambda_2')d_2  = \lambda_0'd_0 +  (\lambda_3'-\lambda_3)d_3.
    \end{equation}
    If $\lambda_0'>0$, then by Lemma \ref{lemma_procedure_1}, we get a contradiction with $[[\lambda_1,\lambda_2,\lambda_3]] \in T$, hence $\lambda_0'=0$. This means that for all $z \in \varphi^{-1}(b)$ we have $0 \notin \text{supp}(z)$, as $(\lambda_0',\lambda_1',\lambda_2', \lambda_3') \in \varphi^{-1}(b)$ was chosen arbitrarily. 

  
    Let $z=(0,\mu_1, \mu_2,  \mu_3) \in \varphi^{-1}(b)$ be a factorization that is not $\mathcal{R}$-related to $z_0$. Repeating the same argument as for $(\lambda_0',\lambda_1',\lambda_2', \lambda_3')$, without loss of generality, we may assume $\lambda_1 \geq \mu_1$, $\lambda_2 \geq \mu_2$, $\lambda_3 < \mu_3$. This forces $\mu_1= \mu_2= \lambda_3 =0$, as otherwise $z \, \mathcal{R} \, z_0$. Hence
    \begin{equation}\label{betticharacter_eq03}
       b= \lambda_1d_1 + \lambda_2d_2   =  \mu_3d_3
    \end{equation}
    where $\lambda_1, \lambda_2, \mu_3> 0$ --- if $\lambda_1 =0$ or $\lambda_2 = 0$, we derive a contradiction as in \eqref{betticharacter_eq00}.
    If $ \mu_3 > a_{33}$, then 
    \begin{equation}
        b  =  a_{30}d_0 + a_{31}d_1 + a_{32}d_2 +  (\mu_3-a_{33})d_3
    \end{equation}
    where $(a_{30},a_{31},a_{32},0) \in \varphi^{-1}(a_{33}d_3)$. We also have $a_{30} = 0$, since $0 \notin \text{supp}(z)$ for any $z \in \varphi^{-1}(b)$. Hence $\max \{a_{31},a_{32}\} > 0$, which implies $z_0 \, \mathcal{R} \, (0,a_{31},a_{32},\mu_3-a_{33}) \, \mathcal{R} \,z$, which is a contradiction. Thus $\mu_3 = a_{33}$, and hence $b = a_{33}d_3 \in U$.
\end{proof}    


\begin{thm}\label{theorembettifind_primary}
If $S$ is primary, then $\textup{Betti}(S) \subseteq U$. 
\end{thm}



\begin{proof}
    Suppose that there exists an element $b \in \text{Betti}(S)$ that is not in $U$.
    There exists a factorization $z_0 =(0,\kappa_1,\kappa_2,\kappa_3) \in \varphi^{-1}(b)$; otherwise, if every factorization $z \in \varphi^{-1}(b)$ had $0 \in \text{supp}(z)$, then all factorizations would be $\mathcal{R}$-related, which cannot happen for a Betti element. By Lemma \ref{lemmabettifind_primary}, we have $[[\kappa_1,\kappa_2,\kappa_3]] \notin T$, hence $[[\kappa_1,\kappa_2,\kappa_3]]$ is in the region associated to some point $P \in V$.

    \begin{itemize}

    

    \item  If $P =(-\lambda_1,\lambda_2,\lambda_3)$, where
    \begin{equation}\label{betticharacter_eq1}
        \lambda_0d_0 + \lambda_1d_1 = \lambda_2d_2 + \lambda_3d_3
    \end{equation}
    with $\lambda_i \in \mathbb{Z}_{>0}$, we get
     \begin{equation}
        b =\lambda_0d_0 + (\lambda_1+\delta_1)d_1 + \delta_2d_2 + \delta_3d_3  =  \delta_1d_1 + (\lambda_2+\delta_2)d_2 + (\lambda_3+\delta_3)d_3
    \end{equation}
    for some $\delta_i \in \mathbb{N}$ with $\max \{\delta_1,\delta_2,\delta_3\} > 0$. The factorizations $z_1 =(\lambda_0,\lambda_1+\delta_1,\delta_2,\delta_3)$ and $z_2 =(0, \delta_1,\lambda_2+\delta_2,\lambda_3+\delta_3)$ have a nonempty intersection of supports, and hence they are $\mathcal{R}$-related. Moreover, we have $0,1,2,3 \in \textup{supp}(z_1) \cup \textup{supp}(z_2)$, so every factorization of $b$ will be $\mathcal{R}$-related to $z_1$ and $z_2$. This contradicts $b \in \textup{Betti}(S)$, as it must have at least two distinct $\mathcal{R}$-classes. 
    
    The cases $P=(\lambda_1,-\lambda_2,\lambda_3)$ and $P=(\lambda_1,\lambda_2,-\lambda_3)$ follow analogously.
    
    \item If $P =(a_{01},a_{02},a_{03})$, where $(0,a_{01},a_{02},a_{03}) \in \varphi^{-1}(a_{00}d_0)$, we get
     \begin{equation}
        b =a_{00}d_0 + \delta_1d_1 + \delta_2d_2 + \delta_3d_3  =  (a_{01}+\delta_1)d_1 + (a_{02}+\delta_2)d_2 + (a_{03}+\delta_3)d_3
    \end{equation}
    for some $\delta_i \in \mathbb{N}$ with $\max \{\delta_1,\delta_2,\delta_3\} > 0$. The factorizations $z_1 =(a_{00},\delta_1,\delta_2,\delta_3)$ and $z_2 =(0,a_{01}+\delta_1,a_{02}+\delta_2,a_{03}+\delta_3)$ have a nonempty intersection of supports, and hence they are $\mathcal{R}$-related. Because $b \in \textup{Betti}(S)$, there exists a factorization $z \in \varphi^{-1}(b)$ which is not $\mathcal{R}$-related to them. 
    
    In this case, as in the proof of Lemma \ref{lemmabettifind_primary}, it will be more convenient to drop the assumption that $S$ is ordered --- we only suppose that $a_{ii}=b_{ii}$ for at least two distinct $i \in \{1,2,3\}$ and that  $T = R \cap \{x \leq a_{11},y\leq a_{22}, z \leq a_{33}\}$. This way, without loss of generality, we can suppose that $\delta_1 > 0$. As $a_{00} >0$, we have $0,1 \notin \text{supp}(z)$. 

    First, suppose that $z = (0,0,\mu_2,\mu_3)$ where $\mu_2,\mu_3>0$. This implies $z_1 = (a_{00},\delta_1,0,0 )$ and $z_2 =(0,a_{01}+\delta_1,0,0)$. 
    If $\mu_2 \geq a_{22}$, then
   \begin{equation}\label{betticharacter_eq40}
         b =\mu_2d_2 + \mu_3d_3 = a_{20}d_0 + a_{21}d_1 +(\mu_2-a_{22})d_2 + (\mu_3+a_{23})d_3
   \end{equation}
   where $(\mu_2-a_{22}) \geq 0$ and $(a_{20},a_{21},0,a_{23}) \in \varphi^{-1}(a_{22}d_2)$.
   If $\max \{a_{20},a_{21} \}>0$, then $z_1 \, \mathcal{R} \,(a_{20},a_{21},$ $ \mu_2-a_{22},\mu_3+a_{23}) \, \mathcal{R} \, z$ (recall $\mu_3>0$), which is a contradiction. Hence $a_{20}=a_{21}=0$ and $a_{22}d_2 = a_{23}d_3$, which implies $a_{23} \geq a_{33}$. Moreover, since $(a_{20},a_{21},0,a_{23})$ was chosen arbitrarily, the conclusion $a_{20}=0$ gives $a_{22} < b_{22}$. Now, we have
   \begin{equation}\label{betticharacter_eq401}
        b=  a_{30}d_0 + a_{31}d_1 + (\mu_2-a_{22} + a_{32})d_2 + (\mu_3+a_{23}-a_{33})d_3
   \end{equation}
   where all coefficients above are nonnegative (since $a_{23} \geq a_{33}$) and $(a_{30},a_{31},a_{32},0) \in \varphi^{-1}(a_{33}d_3)$. As in \eqref{betticharacter_eq40}, this yields $a_{30} = 0$ and hence $a_{33} < b_{33}$. Since we also have $a_{22} < b_{22}$, we get a contradiction with the hypothesis.
   
   If $\mu_3 \geq a_{33}$ we proceed analogously and get a contradiction, thus $\mu_2 < a_{22}$ and $ \mu_3 < a_{33}$. 
   Consider the cube $[[0,\mu_2,\mu_3]]$. If it is in $T$, then $b \in U$ by Lemma \ref{lemmabettifind_primary}. Hence $[[0,\mu_2,\mu_3]]$ is in the region associated to a point $Q \in V$.
   As $\mu_2 < a_{22}$ and $ \mu_3 < a_{33}$, we have $Q =(-\nu_1,\nu_2,\nu_3)$, where
    \begin{equation}\label{betticharacter_eq41}
        - \nu_1d_1 + \nu_2d_2 +  \nu_3d_3   =  \nu_0d_0
    \end{equation}
    $\nu_0,\nu_2,\nu_3 \in \mathbb{Z}_{>0}$, $\nu_1 \in \mathbb{N}$ (possibly $\nu_1=0$), and $\mu_2 \geq \nu_2$, $\mu_3 \geq \nu_3$.  We can write 
    \begin{equation}\label{betticharacter_eq42}
        b =  \nu_0d_0 + \nu_1d_1 + (\mu_2-\nu_2)d_2 + (\mu_3-\nu_3)d_3.
    \end{equation}
    We have $\max \{\mu_2-\nu_2,\mu_3-\nu_3 \}>0$, as $b \notin U$, hence $z \, \mathcal{R} \, (\nu_0,\nu_1,\mu_2-\nu_2,\mu_3-\nu_3) \, \mathcal{R}  \, z_1$, which is a contradiction. 
    

    

    
    
   
    Next, suppose that $z = (0,0,\mu_2,0)$. 
    If $\mu_2 = a_{22} $, then $b = a_{22}d_2 \in U$, so $\mu_2 > a_{22}$. We can write
    \begin{equation}\label{betticharacter_eq6}
         b=\mu_2d_2 = a_{20}d_0 + a_{21}d_1 + (\mu_2-a_{22})d_2 + a_{23}d_3.
    \end{equation}
    where $(a_{20},a_{21},0,a_{23}) \in \varphi^{-1}(a_{22}d_2)$.
    Because $(\mu_2-a_{22}) > 0$, we have $z \, \mathcal{R} \, (a_{20},a_{21},\mu_2-a_{22},a_{23})$. Since $z$ is not $\mathcal{R}$-related to $z_1$, we have $a_{20} = a_{21} = 0$ (recall $a_{00},\delta_1 >0$). This gives the factorization $(0,0,\mu_2-a_{22},a_{23}) \in \varphi^{-1}(b)$ that has $(\mu_2-a_{22}),a_{23} > 0$ and which is not $\mathcal{R}$-related to $z_1$ nor $z_2$, as it is $\mathcal{R}$-related to $z$. Thus we reduce to the previously treated case $z =(0,0,\mu_2,\mu_3)$. 

    The case $z = (0,0,0,\mu_3)$ follows analogously. 
    

    \item   If $P =(b_{11},-b_{12},-b_{13})$, where $(b_{10},0,b_{12},b_{13}) \in \varphi^{-1}(b_{11}d_1)$ with $b_{10}>0$, we get
    \begin{equation}
        b = (b_{11}+\delta_1)d_1 +   \delta_2d_2 + \delta_3d_3  =  b_{10}d_0 + \delta_1d_1+ (b_{12}+\delta_2)d_2 + (b_{13}+\delta_3)d_3
    \end{equation}
    for some $\delta_i \in \mathbb{N}$ with $\max \{\delta_1,\delta_2,\delta_3\} > 0$. The factorizations $z_1 =(0,b_{11} +\delta_1,\delta_2,\delta_3)$ and $z_2 =(b_{10},\delta_1,b_{12}+\delta_2,b_{13}+\delta_3)$ have a nonempty intersection of supports, and hence they are $\mathcal{R}$-related. Because $b \in \text{Betti}(S)$, there exists a factorization $z \in \varphi^{-1}(b)$ that is not $\mathcal{R}$-related to them. 
        
        If $b_{12},b_{13}>0$, then $0,1,2,3 \in \textup{supp}(z_1) \cup \textup{supp}(z_2)$, which implies that $b$ has only one $\mathcal{R}$-class.
        
        If $b_{12}>0,b_{13} = 0$, then $0,1,2 \in \textup{supp}(z_1) \cup \textup{supp}(z_2)$ and hence $z = (0,0,0,\mu_3)$. 
        If $\mu_3 = a_{33}$, then $b = a_{33}d_3 \in U$, hence $\mu_3 > a_{33}$. Now, for $(a_{30},a_{31},a_{32},0) \in \varphi^{-1}(a_{33}d_3)$, we have $z \, \mathcal{R} \, (a_{30},a_{31},a_{32},\mu_3-a_{33}) \, \mathcal{R} \, z_1  \, \mathcal{R} \, z_2$, since $\max\{a_{30},a_{31},a_{32}\}>0$, which is a contradiction.

        If $b_{12}=0, b_{13} > 0$, we proceed analogously as above.
        
        
        If $b_{12} = b_{13} = 0$, then $P = (b_{11},0,0)$ and $b_{10} \geq a_{00}$. If $b_{10} = a_{00}$, then we are in the case $P = (a_{01},a_{02},a_{03})$ and we can proceed analogously as there. If $b_{10} > a_{00}$, then $(b_{10}-a_{00},a_{01},a_{02},a_{03}) \in \varphi^{-1}(b_{11}d_1)$ with $\max \{a_{01},a_{02},a_{03}\}>0$. If $a_{01} > 0$, we get a contradiction with the minimality of $b_{11}$ (since $(b_{10}-a_{00})>0$). Hence $(b_{10}-a_{00},0,a_{02},a_{03}) \in \varphi^{-1}(b_{11}d_1)$ where $(b_{10}-a_{00})>0$ and $\max \{a_{02},a_{03}\} > 0$. Therefore we can replace $P = (b_{11},0,0)$ with $(b_{11},-a_{02},-a_{03})$ and proceed as in the previous cases of this bullet-point.

        The case $P = (-b_{21},b_{22},-b_{23})$ follows analogously.
        
        

    \item   Suppose $P =(-a_{31},-a_{32},a_{33})$, where $(a_{30},a_{31},a_{32},0) \in \varphi^{-1}(a_{33}d_3)$. If $a_{33}=b_{33}$, then we can proceed as in the previous bullet-point. Thus, suppose $a_{33}<b_{33}$, which implies $a_{30} = 0$. We get
    \begin{equation}
        b = \delta_1d_1 +   \delta_2d_2 + (a_{33} +\delta_3)d_3  =  (a_{31}+\delta_1)d_1+ (a_{32}+\delta_2)d_2 + \delta_3d_3
    \end{equation}
    for some $\delta_i \in \mathbb{N}$ with $\max \{\delta_1,\delta_2,\delta_3\} > 0$. The factorizations $z_1 =(0,\delta_1,\delta_2,a_{33} +\delta_3)$ and $z_2 =(0,a_{31}+\delta_1,a_{32}+\delta_2,\delta_3)$ have a nonempty intersection of supports, and hence they are $\mathcal{R}$-related. Because $b \in \text{Betti}(S)$, there exists a factorization $z \in \varphi^{-1}(b)$ that is not $\mathcal{R}$-related to them. 
        
    If $a_{31},a_{32}>0$, then $1,2,3 \in \textup{supp}(z_1) \cup \textup{supp}(z_2)$, hence $z = (\mu_0,0,0,0)$. Consider the cube $[[a_{31}+\delta_1,a_{32}+\delta_2,\delta_3]]$ which has label $b$. By Lemma \ref{lemma_catenary_help_primary}, we have $[[a_{31}+\delta_1,a_{32}+\delta_2,\delta_3]] \notin T$. If $\delta_3 < a_{33}$, then $[[a_{31}+\delta_1,a_{32}+\delta_2,\delta_3]]$ lies in the region associated to some point $Q \in V$, $Q \neq (-a_{31},-a_{32},a_{33})$, and we can proceed as in one of the cases covered so far. If $\delta_3 \geq a_{33}$, then there exists $\lambda \in \mathbb{Z}_{>0}$ such that $a_{33} >\delta_3 - \lambda a_{33} \geq 0$. Now, the cube $[[ (\lambda+1)a_{31}+\delta_1,(\lambda+1)a_{32}+\delta_2,\delta_3 - \lambda a_{33}]]$ has label $b$ and is in the region associated to some point $Q \in V$, $Q \neq (-a_{31},-a_{32},a_{33})$, hence we can proceed as in one of the cases covered so far.

    If $a_{31} = 0$, then $a_{33}d_3 = a_{32}d_2$, which implies $a_{32} \geq a_{22} = b_{22}$. We get 
    \begin{equation}
        a_{33}d_3 = b_{20}d_0 + b_{21}d_1 + (a_{32}-b_{22})d_2 + b_{23}d_3
    \end{equation}
    where $(b_{20},b_{21},0,b_{23}) \in \varphi^{-1}(b_{22}d_2)$ with $b_{20} > 0$, which contradicts $a_{33} < b_{33}$.

    The case $a_{32}=0$ follows analogously. \qedhere


\end{itemize}
\end{proof}

\begin{lemma}\label{lemma_catenary_help_primary}
    Suppose that $S$ is primary.  If $(-\lambda_1,\lambda_2,\lambda_3) \in V$ satisfies
    $$ \lambda_0d_0 + \lambda_1d_1 = \lambda_2d_2 + \lambda_3d_3$$
    with $\lambda_i \in \mathbb{Z}_{>0}$ and $\lambda_0 < a_{00}$, then for any $(\mu_0,\mu_1,\mu_2,\mu_3) \in \varphi^{-1}(\lambda_0d_0 + \lambda_1d_1)$ we have $(\mu_0,\mu_1,\mu_2,\mu_3) = (0,0,\lambda_2,\lambda_3)$ or $\mu_2=\mu_3 = 0$. An analogous statement holds for $(\lambda_1,-\lambda_2,\lambda_3)$ and $(\lambda_1,\lambda_2,-\lambda_3) \in V$.
\end{lemma}  


\begin{proof}
Note that $\lambda_2< a_{22}$ and $\lambda_3 < a_{33}$ by the definition of $V$. 

Suppose $\mu_0>0$. If $\mu_2 > 0$, then by Lemma \ref{lemma_procedure_1}, no element of the Ap{\'e}ry set is a label in the region associated to $(-\mu_1,\lambda_2-\mu_2,\lambda_3-\mu_3)$. This contradicts the fact $[[0,\lambda_2-1, \lambda_3]] \in T$, which holds since $(-\lambda_1,\lambda_2,\lambda_3) \in V$. If $\mu_3 > 0$, we get an analogous contradiction, hence $\mu_2=\mu_3 = 0$.

Thus, suppose $\mu_0 = 0$. If $\mu_2 \geq \lambda_2$, then
\begin{equation}
     (\lambda_3-\mu_3)d_3=\mu_1d_1 + (\mu_2 - \lambda_2)d_2
\end{equation}
where all coefficients above are nonnegative. If $(\lambda_3-\mu_3) > 0$, then we get a contradiction with $\lambda_3 < a_{33}$. Hence $(\lambda_3-\mu_3) = 0$, which gives $(\mu_0,\mu_1,\mu_2,\mu_3) = (0,0,\lambda_2,\lambda_3)$. If $\mu_3 \geq \lambda_3$, we get the same thing; therefore, suppose that $\mu_2 < \lambda_2$ and $\mu_3 < \lambda_3$. We have
\begin{equation}
   \lambda_0d_0 + (\lambda_1-\mu_1)d_1 = \mu_2d_2 + \mu_3d_3.
\end{equation}
If $(\lambda_1-\mu_1) \leq 0$, then we get a contradiction with $\lambda_0 < a_{00}$, hence $(\lambda_1-\mu_1) > 0$. Since $\lambda_0>0$, by Lemma \ref{lemma_procedure_1}, no cube labeled with an element of the Ap{\'e}ry set is in the region associated to $(\mu_1 - \lambda_1,\mu_2,\mu_3)$. This contradicts the fact $[[0,\lambda_2-1,\lambda_3-1]] \in T$, which holds since $(-\lambda_1,\lambda_2,\lambda_3) \in V$.

The same argument applies to $(\lambda_1,-\lambda_2,\lambda_3)$ and $(\lambda_1,\lambda_2,-\lambda_3) \in V$.
\end{proof}




\begin{lemma}\label{lemma_catenary_primary}
    Suppose that $S$ is primary. If $(-\lambda_1,\lambda_2,\lambda_3) \in V$ satisfies
    $$ \lambda_0d_0 + \lambda_1d_1 = \lambda_2d_2 + \lambda_3d_3$$
    with $\lambda_i \in \mathbb{Z}_{>0}$, $\lambda_0 < a_{00}$, and $ \lambda_1 < a_{11}$, then $\varphi^{-1}(\lambda_0d_0 + \lambda_1d_1) = \{(\lambda_0,\lambda_1,0,0), (0,0,\lambda_2,\lambda_3) \} $. An analogous statement holds for $(\lambda_1,-\lambda_2,\lambda_3)$ and $(\lambda_1,\lambda_2,-\lambda_3) \in V$. 
\end{lemma}

\begin{proof}
Note that $\lambda_2< a_{22}$ and $\lambda_3 < a_{33}$ by the definition of $V$.
Suppose that there exists a factorization $ (\mu_0,\mu_1,\mu_2,\mu_3) \in \varphi^{-1}(\lambda_0d_0 + \lambda_1d_1)$ distinct from the two above. 
By Lemma \ref{lemma_catenary_help_primary}, we have $\mu_2 = \mu_3 = 0 $. 
If $\lambda_0 \geq \mu_0$, then
\begin{equation}
    (\lambda_0- \mu_0) d_0 =  (\mu_1 - \lambda_1)d_1
\end{equation}
where all coefficients above are nonnegative. Indeed, they are positive, since otherwise $(\mu_0,\mu_1,\mu_2,\mu_3) = (\lambda_0,\lambda_1,0,0)$. This contradicts $\lambda_0 < a_{00}$. If $\lambda_0 < \mu_0$, then $\lambda_1 > \mu_1$, and we get a contradiction with $\lambda_1 < a_{11}$. 

The same argument applies to $(\lambda_1,-\lambda_2,\lambda_3)$ and $(\lambda_1,\lambda_2,-\lambda_3) \in V$.
\end{proof}


\begin{lemma}\label{lemma_catenarycase1_primary}
    Suppose that $S$ is primary. Consider $(-\lambda_1,\lambda_2,\lambda_3) \in V$ satisfying
    $$ \lambda_0d_0 + \lambda_1d_1 = \lambda_2d_2 + \lambda_3d_3$$
    with $\lambda_i \in \mathbb{Z}_{>0}$, and $ \lambda_0 \geq a_{00}$ or $ \lambda_1 \geq a_{11}$. If $a_{00}d_0 \neq a_{11}d_1$, then $\lambda_0d_0 + \lambda_1d_1 \notin \textup{Betti}(S)$. Otherwise, if $a_{00}d_0 = a_{11}d_1$, then $\lambda_0d_0 + \lambda_1d_1 \in \textup{Betti}(S)$ and $\varphi^{-1}(\lambda_0d_0 + \lambda_1d_1) = \mathcal{F} \cup \{(0,0,\lambda_2,\lambda_3)\}$, where 
    $$\mathcal{F} =   \{(\lambda_0 + \lambda a_{00},\lambda_1 - \lambda a_{11},0,0) \mid \lambda \in \mathbb{Z}, \lambda_0 + \lambda a_{00} \geq 0, \lambda_1 - \lambda a_{11} \geq 0 \}. $$ 
    An analogous statement holds for $(\lambda_1,-\lambda_2,\lambda_3)$ and $(\lambda_1,\lambda_2,-\lambda_3) \in V$.
\end{lemma}


\begin{proof}
    Note that $\lambda_2< a_{22}$ and $\lambda_3 < a_{33}$ by the definition of $V$.
    
    Suppose $a_{00}d_0 \neq a_{11}d_1$ and $\lambda_0 \geq a_{00}$. Consider $(0,a_{01},a_{02},a_{03}) \in \varphi^{-1}(a_{00}d_0)$. If $\max \{a_{02},a_{03}\} > 0$, then
    \begin{equation}
        (\lambda_0,\lambda_1,0,0) \, \mathcal{R} \, (\lambda_0 - a_{00},\lambda_1 + a_{01},a_{02},a_{03}) \, \mathcal{R} \, (0,0,\lambda_2,\lambda_3)
    \end{equation}
    which results in only one $\mathcal{R}$-class. Hence $a_{02}=a_{03}=0$ and $a_{00}d_0 = a_{01}d_1$. Since $a_{00}d_0 \neq a_{11}d_1$, this gives $a_{01} > a_{11}$. Consider $(a_{10},0,a_{12},a_{13}) \in \varphi^{-1}(a_{11}d_1)$. If $a_{10}>0$, then we get a contradiction with the minimality of $a_{00}$, hence $a_{10}=0$. This implies  $\max \{a_{12},a_{13}\} > 0$ and hence 
    \begin{equation}
        (\lambda_0,\lambda_1,0,0) \, \mathcal{R} \, (\lambda_0 - a_{00},\lambda_1 + a_{01},0,0) \, \mathcal{R} \, (\lambda_0 - a_{00},\lambda_1 + a_{01}-a_{11},a_{12},a_{13}) \, \mathcal{R} \, (0,0,\lambda_2,\lambda_3)
    \end{equation}
    which results in only one $\mathcal{R}$-class. 
    
    If $a_{00}d_0 \neq a_{11}d_1$ and $\lambda_0 < a_{00}$, then $\lambda_1 \geq a_{11}$ and we can apply the same argument as above. Thus the first result follows.

    Suppose $a_{00}d_0 = a_{11}d_1$. Clearly $\mathcal{F} \subseteq \varphi^{-1}(\lambda_0d_0 + \lambda_1d_1)$. Suppose that there exists $(\mu_0,\mu_1,\mu_2,\mu_3) \in \varphi^{-1}(\lambda_0d_0 + \lambda_1d_1)$ that is not in $\mathcal{F} \cup \{(0,0,\lambda_2,\lambda_3)\}$.  By Lemma \ref{lemma_catenary_help_primary}, we get $\mu_2=\mu_3 = 0$ --- we can apply this lemma, because in $\mathcal{F}$, there exists a factorization that satisfies its assumptions. Then, as $(\mu_0,\mu_1,0,0) \notin \mathcal{F}$, we can find $\lambda \in \mathbb{Z}$ such that $\lambda_0 + (\lambda+1)a_{00} > \mu_0 > \lambda_0 + \lambda a_{00}$. This gives
    \begin{equation}
        \lambda_0d_0 + \lambda_1d_1 = (\mu_0 - \lambda a_{00})d_0 + (\mu_1 + \lambda a_{11})d_1 \implies (\mu_0 -\lambda a_{00} - \lambda_0)d_0 = (\lambda_1 - \mu_1 - \lambda a_{11})d_1  
    \end{equation}
    which gives a contradiction, since $ a_{00} > (\mu_0 - \lambda_0 - \lambda a_{00}) > 0$.

    The same argument applies to $(\lambda_1,-\lambda_2,\lambda_3)$ and $(\lambda_1,\lambda_2,-\lambda_3) \in V$.
\end{proof}

Lemma \ref{lemma_catenary_minrelations} gives $a_{ii}d_i \in \textup{Betti}(S)$ for $i = 0,1,2,3$. Combined with Lemmas \ref{lemma_catenary_primary} and \ref{lemma_catenarycase1_primary}, we know exactly which elements of $U$ are Betti elements. Moreover, we get that every Betti element that corresponds to a $(2,2)$-type point has exactly two $\mathcal{R}$-classes.



\begin{lemma}\label{lemma_catenaryT_primary}
    Suppose that $S$ is primary. If $(-\lambda_1,\lambda_2,\lambda_3) \in V$ satisfies
    $$ \lambda_0d_0 + \lambda_1d_1 = \lambda_2d_2 + \lambda_3d_3 \notin \textup{Betti}(S)$$
    with $\lambda_i \in \mathbb{Z}_{>0}$, then $\lambda_0 = a_{00}$ and $\lambda_1 < a_{11}$. An analogous statement holds for $(\lambda_1,-\lambda_2,\lambda_3)$ and $(\lambda_1,\lambda_2,-\lambda_3) \in V$. 
\end{lemma}

\begin{proof}

The condition $\lambda_0d_0 + \lambda_1d_1 \notin \textup{Betti}(S)$ implies $a_{00}d_0 \neq a_{11}d_1$, and $\lambda_0 \geq a_{00}$ or $\lambda_1 \geq a_{11}$, by Lemmas \ref{lemma_catenary_primary} and \ref{lemma_catenarycase1_primary}. Also, note that $\lambda_2< a_{22} $ and $\lambda_3 < a_{33} $ by the definition of $V$.

Suppose $\lambda_1 \geq a_{11}$. For $(a_{10},0,a_{12},a_{13}) \in \varphi^{-1}(a_{11}d_1)$, we have
\begin{equation}
     (\lambda_0+a_{10})d_0 + (\lambda_1-a_{11})d_1 =  (\lambda_2-a_{12})d_2 + (\lambda_3-a_{13})d_3
\end{equation}
where $(\lambda_1-a_{11}) \geq 0$. 
If $(\lambda_2-a_{12}) < 0$, then $\lambda_3 \geq (\lambda_3-a_{13}) \geq a_{33}$, which is a contradiction. If $(\lambda_3-a_{13}) < 0$, we get the same contradiction. Hence $(\lambda_2-a_{12}),(\lambda_3-a_{13}) \geq 0$.
If $\max \{a_{12},a_{13}\} > 0$, then by Lemma \ref{lemma_procedure_1}, no cube labeled with an element of the Ap{\'e}ry set is in the region associated to $(-\lambda_1+a_{11},\lambda_2-a_{12},\lambda_3-a_{13})$. This contradicts the fact $[[0,\lambda_2-1,\lambda_3-1]] \in T$, which holds since $(-\lambda_1,\lambda_2,\lambda_3) \in V$. 
Hence $a_{12} =a_{13}= 0$ and $a_{11}d_1 = a_{10}d_0$. This gives $a_{10} > a_{00}$, as $a_{00}d_0 \neq a_{11}d_1$. For $(0,a_{01},a_{02},a_{03}) \in \varphi^{-1}(a_{00}d_0)$, we have
\begin{equation}
    a_{11}d_1 = (a_{10}-a_{00})d_0 + a_{01}d_1 + a_{02}d_2 + a_{03}d_3
\end{equation}
which gives $a_{01} = 0$, as otherwise we would get a contradiction with the minimality of $a_{11}$. Therefore, we have $\max \{ a_{02},a_{03}\} > 0$ and
\begin{equation}
    (\lambda_0+a_{10}-a_{00})d_0 + (\lambda_1-a_{11} + a_{01})d_1 =  (\lambda_2-a_{02})d_2 + (\lambda_3-a_{03})d_3.
\end{equation}
Since $(\lambda_0+a_{10}-a_{00}) > 0$, no cube labeled with an element of the Ap{\'e}ry set is in the region associated to $(-\lambda_1+a_{11} - a_{01},\lambda_2-a_{02},\lambda_3-a_{03})$, by Lemma \ref{lemma_procedure_1}. Again, this contradicts the fact $[[0,\lambda_2-1,\lambda_3-1]] \in T$. Therefore $\lambda_1 <a_{11}$. 

If $\lambda_0 > a_{00}$, then for $(0,a_{01},a_{02},a_{03}) \in \varphi^{-1}(a_{00}d_0)$, we have
\begin{equation}
     (\lambda_0-a_{00})d_0 + (\lambda_1+a_{01})d_1 =  (\lambda_2-a_{02})d_2 + (\lambda_3-a_{03})d_3
\end{equation} 
where $(\lambda_0-a_{00}) > 0$. Proceeding as in the previous case, we obtain $a_{02}=a_{03}=0$ and $a_{01} > a_{11}$. Then, following the same approach further, we get $\max \{a_{12},a_{13}\} > 0$ and at last a contradiction. Therefore $\lambda_0 = a_{00}$.

The same argument applies to $(\lambda_1,-\lambda_2,\lambda_3)$ and $(\lambda_1,\lambda_2,-\lambda_3) \in V$.
\end{proof}


\subsection{Secondary case}


\begin{lemma}\label{lemmabettifind_secondary}
   Suppose that $S$ satisfies $a_{11} = b_{11}$ and $a_{22}d_2 = a_{33}d_3$. If $b \in \textup{Betti}(S)$ labels a cube in $R$, then $b = a_{22}d_2 = a_{33}d_3$.
\end{lemma}
    
\begin{proof}
    Consider a factorization $(0,\lambda_1,\lambda_2, \lambda_3) \in \varphi^{-1}(b)$ with $[[\lambda_1,\lambda_2,\lambda_3]] \in R$. Let $(\lambda_0',\lambda_1',\lambda_2', \lambda_3')$ be a distinct factorization of $b$. We have
    \begin{equation}\label{betticharacter_eq0a}
        (\lambda_1-\lambda_1')d_1 +  (\lambda_2-\lambda_2')d_2 +  (\lambda_3-\lambda_3')d_3 = \lambda_0'd_0.
    \end{equation}
    \begin{itemize}
    \item 
    If all coefficients on the left-hand side of \eqref{betticharacter_eq0a} are nonnegative, then $\lambda_0'>0$, as the two factorizations are distinct. This contradicts $[[\lambda_1,\lambda_2,\lambda_3]] \in R$ by Lemma \ref{lemma_procedure_1}.

    \item 
    Suppose that exactly one coefficient on the left-hand side of \eqref{betticharacter_eq0a} is nonnegative. If $(\lambda_1-\lambda_1')$ is the nonnegative coefficient, then $\lambda_1 \geq  (\lambda_1-\lambda_1') \geq a_{11}=b_{11}$, which contradicts $[[\lambda_1,\lambda_2,\lambda_3]] \in R$. Therefore, without loss of generality, suppose that $(\lambda_2-\lambda_2')$ is the nonnegative coefficient. This gives 
    \begin{equation}
          (\lambda_2-\lambda_2')d_2    = \lambda_0'd_0 + (\lambda_1'-\lambda_1)d_1 + (\lambda_3'-\lambda_3)d_3
    \end{equation}
    where every coefficient above is nonnegative, which implies $ (\lambda_2-\lambda_2') \geq a_{22}$ (since $(\lambda_2-\lambda_2') \neq 0$, as $(\lambda_3'-\lambda_3) > 0$). If $\lambda_0' > 0$, then $\lambda_2 \geq  (\lambda_2-\lambda_2') \geq b_{22}$, which contradicts $[[\lambda_1,\lambda_2,\lambda_3]] \in R$, hence $\lambda_0' = 0$.  
    There exists $\lambda \in \mathbb{N}$ such that $a_{22} > \lambda_2-\lambda_2' - \lambda a_{22} \geq 0$. Using $a_{22}d_2 = a_{33}d_3$, we obtain 
    \begin{equation}
          (\lambda_2-\lambda_2'- \lambda a_{22} )d_2    =  (\lambda_1'-\lambda_1)d_1 + (\lambda_3'-\lambda_3-\lambda a_{33} )d_3.
    \end{equation}
    If $(\lambda_3'-\lambda_3-\lambda a_{33} ) > 0$, then we get a contradiction with $a_{22} > \lambda_2-\lambda_2' - \lambda a_{22}$, hence every coefficient below is nonnegative.
    \begin{equation}
          (\lambda_2-\lambda_2'- \lambda a_{22} )d_2 + ( \lambda_3 -\lambda_3'+ \lambda a_{33})d_3    =  (\lambda_1'-\lambda_1)d_1.
    \end{equation}
    If $(\lambda_1'-\lambda_1) >0$, then $\lambda_1' \geq (\lambda_1'-\lambda_1) \geq a_{11} = b_{11}$. By Theorem \ref{theoremprocedure}, $R$ is the collection of all cubes labeled with an element of the Ap{\'e}ry set. This implies that $b$ is an element of the Ap{\'e}ry set, as it labels $[[\lambda_1,\lambda_2,\lambda_3]] \in R$. However, $b$ also labels $[[\lambda_1',\lambda_2',\lambda_3']]$ (recall $\lambda_0' = 0$), which is not in $R$ since $\lambda_1' \geq b_{11}$ --- contradiction.

    Hence $(\lambda_1'-\lambda_1) = 0$ and $(\lambda_0',\lambda_1',\lambda_2', \lambda_3') = (0,\lambda_1,\lambda_2 - \lambda a_{22} , \lambda_3 + \lambda a_{33} )$. As $(\lambda_0',\lambda_1',\lambda_2', \lambda_3') \in \varphi^{-1}(b)$ was chosen arbitrarily, we obtain
    \begin{equation}\label{betticharacter_factor}
        \varphi^{-1}(b) = \{ (0,\lambda_1,\lambda_2 - \lambda a_{22} , \lambda_3 + \lambda a_{33} ) \mid \lambda \in \mathbb{Z}, \lambda_2 - \lambda a_{22} \geq 0 , \lambda_3 + \lambda a_{33} \geq 0\}.
    \end{equation}
    Therefore $b = a_{22}d_2 = a_{33}d_3$, as otherwise $b$ would have only one $\mathcal{R}$-class, which would contradict $b \in \textup{Betti}(S)$.

    \item 
    Suppose that exactly two coefficients on the left-hand side of \eqref{betticharacter_eq0a} are nonnegative. 
    
    If $(\lambda_1-\lambda_1') < 0$, then
    \begin{equation}
        (\lambda_2-\lambda_2')d_2 +  (\lambda_3-\lambda_3')d_3 = \lambda_0'd_0 +  (\lambda_1'-\lambda_1)d_1
    \end{equation}
    where every coefficient above is nonnegative. If $\lambda_0' > 0$, we get a contradiction with $[[\lambda_1,\lambda_2,\lambda_3]] \in  R$ by Lemma \ref{lemma_procedure_1}. Hence $\lambda_0' = 0$, which implies $\lambda_1' \geq (\lambda_1'-\lambda_1) \geq a_{11} = b_{11}$. As before, $b$ labels $[[\lambda_1',\lambda_2',\lambda_3']]$, which is not in $R$ since $\lambda_1' \geq b_{11}$, contradicting Theorem \ref{theoremprocedure}.  

    Suppose $(\lambda_2-\lambda_2')<0$ (the case $(\lambda_3-\lambda_3') < 0$ follows analogously). This gives
    \begin{equation}
        (\lambda_1-\lambda_1')d_1 +  (\lambda_3-\lambda_3')d_3 = \lambda_0'd_0 +  (\lambda_2'-\lambda_2)d_2
    \end{equation}
    where every coefficient above is nonnegative. If $\lambda_0' > 0$, we get a contradiction with $[[\lambda_1,\lambda_2,\lambda_3]] \in R$ by Lemma \ref{lemma_procedure_1}. Hence $\lambda_0' = 0$, which implies $\lambda_2' \geq (\lambda_2'-\lambda_2) \geq a_{22}$.  There exists $\lambda \in \mathbb{N}$ such that $a_{22} > \lambda_2'-\lambda_2 - \lambda a_{22} \geq 0$. Using $a_{22}d_2 = a_{33}d_3$, we obtain 
    \begin{equation}
               (\lambda_1-\lambda_1')d_1 + (\lambda_3-\lambda_3'-\lambda a_{33} )d_3 = (\lambda_2'-\lambda_2- \lambda a_{22} )d_2.
    \end{equation}
    If $(\lambda_3-\lambda_3'-\lambda a_{33} ) > 0$, then we get a contradiction with $a_{22} > \lambda_2'-\lambda_2 - \lambda a_{22}$, hence every coefficient below is nonnegative.
    \begin{equation}
               (\lambda_1-\lambda_1')d_1 = (\lambda_3'-\lambda_3+\lambda a_{33} )d_3 + (\lambda_2'-\lambda_2- \lambda a_{22} )d_2.
    \end{equation}
     If $(\lambda_1-\lambda_1')$ is positive, then we get $\lambda_1 \geq (\lambda_1-\lambda_1') \geq a_{11} = b_{11}$, which contradicts $[[\lambda_1,\lambda_2,\lambda_3]] \in R$. 
     Hence $(\lambda_1-\lambda_1') = 0$, $(\lambda_0',\lambda_1',\lambda_2', \lambda_3') = (0,\lambda_1,\lambda_2 + \lambda a_{22} , \lambda_3 - \lambda a_{33} )$, and we obtain \eqref{betticharacter_factor}. This again gives $b =a_{22}d_2 =a_{33}d_3$. \qedhere   
    \end{itemize}  
\end{proof}

Let $V' \subset \mathbb{Z}^3$ be a minimal set (with respect to inclusion) of points $(x,y,z)$ such that $xd_1 + yd_2 +zd_3$ is a positive multiple of $d_0$ and such that removing from the initial collection of cubes the regions associated to the points of $V'$ yields $R$. 
Define
$$ U' = \{ \max\{ x,0\} d_1 + \max\{y,0\}d_2 + \max\{z,0\}d_3 \mid (x,y,z) \in V' \}.$$
Note that $U'$ is independent of the choice of $V'$. Similarly as for $V$, assume that every point $(x,y,z) \in V'$ such that
$$ \max\{ x,0\} d_1 + \max\{y,0\}d_2 + \max\{z,0\}d_3 \notin \{(a_{00}d_0,b_{11}d_1,b_{22}d_2,b_{33}d_3)\}$$
is a $(2,2)$-type point. Such $V'$ exists by Theorem \ref{theoremprocedure}.

\begin{lemma}\label{lemmaU_secondary}
    Suppose that $S$ satisfies $a_{11} = b_{11}$ and $a_{22}d_2 = a_{33}d_3$. Then $ U' \cup \{a_{33}d_3\} = U \cup \{b_{33}d_3\}$.
\end{lemma}

\begin{proof}
    Recall that the boundary values of $R$ and $T$ in the $x$-, $y$-, and $z$-directions are $b_{11}$, $b_{22}$, $b_{33}$ and $b_{11}$, $b_{22}$, $a_{33}$, respectively. 
    

    Consider $(-\lambda_1,\lambda_2,\lambda_3) \in V'$ where $\lambda_i > 0$ and $\lambda_3 \geq a_{33}$. There exists $\lambda \in \mathbb{N}$ such that $a_{33} > \lambda_3 - \lambda a_{33} \geq 0$. Then, we have $(-\lambda_1,\lambda_2 + \lambda a_{22},\lambda_3-\lambda a_{33}) \in V$, since $T = R \cap \{z \leq a_{33}\}$ and $[[a,b,c]] \in R \cap \{y \geq a_{22}\}$ if and only if $[[a,b-a_{22},c+a_{33}]] \in R \cap \{z \geq a_{33}\}$ by Proposition \ref{prop_secondary}.
    
    We proceed analogously for any $(\lambda_1,-\lambda_2,\lambda_3) \in V'$ with $\lambda_3 \geq a_{33}$  (and $\lambda_i > 0$). We also have that any $(\lambda_1,\lambda_2,-\lambda_3) \in V'$ ($\lambda_i > 0$) is in $V$, since $T = R \cap \{z \leq a_{33}\}$. Consequently, $ U' \cup \{a_{33}d_3\} = U \cup \{b_{33}d_3\}$.    
\end{proof}

\begin{thm}\label{theorembettifind_secondary}
If $S$ is secondary, then $\textup{Betti}(S) \subseteq U \cup \{b_{33}d_3\}$. 
\end{thm}


\begin{proof}
    Suppose that there exists an element $b \in \text{Betti}(S)$ that is not in $U\cup \{b_{33}d_3\}$, which by Lemma \ref{lemmaU_secondary} equals $U' \cup \{a_{33}d_3\}$.
    There exists a factorization $(0,\kappa_1,\kappa_2,\kappa_3) \in \varphi^{-1}(b)$; otherwise, if every factorization $z \in \varphi^{-1}(b)$ had $0 \in \text{supp}(z)$, then all factorizations would be $\mathcal{R}$-related, which cannot happen for a Betti element. By Lemma \ref{lemmabettifind_secondary}, we have $[[\kappa_1,\kappa_2,\kappa_3]] \notin R$; therefore $[[\kappa_1,\kappa_2,\kappa_3]]$ is in the region associated to some point $P \in V'$.  
\begin{itemize}
    \item 
    If $P =(-\lambda_1,\lambda_2,\lambda_3)$, where
    \begin{equation}
        \lambda_0d_0 + \lambda_1d_1 = \lambda_2d_2 + \lambda_3d_3
    \end{equation}
    with $\lambda_i \in \mathbb{Z}_{>0}$, we get
     \begin{equation}
        b =\lambda_0d_0 + (\lambda_1+\delta_1)d_1 + \delta_2d_2 + \delta_3d_3  =  \delta_1d_1 + (\lambda_2+\delta_2)d_2 + (\lambda_3+\delta_3)d_3
    \end{equation}
    for some $\delta_i \in \mathbb{N}$ with $\max \{\delta_1,\delta_2,\delta_3\} > 0$. The factorizations $z_1 =(\lambda_0,\lambda_1+\delta_1,\delta_2,\delta_3)$ and $z_2=(0, \delta_1,\lambda_2+\delta_2,\lambda_3+\delta_3)$ have a nonempty intersection of supports, and hence they are $\mathcal{R}$-related. Moreover, we have $0,1,2,3 \in \textup{supp}(z_1) \cup \textup{supp}(z_2)$, so every factorization of $b$ will be $\mathcal{R}$-related to $z_1$ and $z_2$. This contradicts $b \in \textup{Betti}(S)$, as it must have at least two distinct $\mathcal{R}$-classes. 
    
    The cases $P=(\lambda_1,-\lambda_2,\lambda_3)$ and $P=(\lambda_1,\lambda_2,-\lambda_3)$ follow analogously.
     

\item
    If $P =(-b_{21},b_{22},-b_{23})$, where $(b_{20},b_{21},0,b_{23}) \in \varphi^{-1} (b_{22}d_2)$ with $b_{20}>0$, we get
    \begin{equation}
        b = \delta_1d_1 +   (b_{22}+\delta_2)d_2 + \delta_3d_3  =  b_{20}d_0 + (b_{21}+\delta_1)d_1+ \delta_2d_2 + (b_{23}+\delta_3)d_3
    \end{equation}
    for some $\delta_i \in \mathbb{N}$ with $\max \{\delta_1,\delta_2,\delta_3\} > 0$. The factorizations $z_1 =(0,\delta_1,b_{22}+\delta_2,\delta_3)$ and $z_2 =(b_{20},b_{21}+\delta_1,\delta_2,b_{23}+\delta_3)$ are $\mathcal{R}$-related. 
    Because $b \in \textup{Betti}(S)$, there exists a factorization $z \in \varphi^{-1}(b)$ that is not $\mathcal{R}$-related to them. Let $z_3 = (0,\delta_1,b_{22}+\delta_2 - a_{22},\delta_3+a_{33}) \in \varphi^{-1}(b)$ (as $a_{22}d_2 = a_{33}d_3$). We have $z_1 \, \mathcal{R} \, z_2 \, \mathcal{R} \,   z_3 $ and $0,2,3 \in \textup{supp}(z_1) \cup \textup{supp}(z_2) \cup \textup{supp}(z_3)$. This implies $z = (0,\mu_1,0,0)$. If $\mu_1 = a_{11} = b_{11}$, then $b = b_{11}d_1 \in U$, hence $\mu_1 > a_{11}$. This gives $z \, \mathcal{R} \, (a_{00},\mu_1-a_{11},0,0) \, \mathcal{R} \, z_2$, which is a contradiction.  
     

    The case $P = (-b_{31},-b_{32},b_{33})$ follows analogously.
\item
    If $P = (a_{01},a_{02},a_{03})$, where $(0,a_{01},a_{02},a_{03}) \in \varphi^{-1}(a_{00}d_0)$, then as $S$ is secondary, we have $(a_{01},a_{02},a_{03}) = (a_{11},0,0)$. Thus this case reduces to the next one.
\item    
    If $P =(b_{11},0,0)=(a_{11},0,0)$, we get
     \begin{equation}
        b =a_{00}d_0 + \delta_1d_1 + \delta_2d_2 + \delta_3d_3  =  (a_{11}+\delta_1)d_1 + \delta_2d_2 + \delta_3d_3
    \end{equation}
    for some $\delta_i \in \mathbb{N}$ with $\max \{\delta_1,\delta_2,\delta_3\} > 0$. The factorizations $z_1 =(a_{00},\delta_1,\delta_2,\delta_3)$ and $z_2 =(0,a_{11}+\delta_1,\delta_2,\delta_3)$ have a nonempty intersection of supports, and hence they are $\mathcal{R}$-related. Because $b \in \textup{Betti}(S)$, there exists a factorization $z \in \varphi^{-1}(b)$ that is not $\mathcal{R}$-related to them. As $0,1 \in \textup{supp}(z_1) \cup \textup{supp}(z_2)$, we have $z = (0,0,\mu_2,\mu_3)$. If $[[0,\mu_2,\mu_3]] \in R$, then $b =a_{22}d_2=a_{33}d_3 \in U$ by Lemma \ref{lemmabettifind_secondary}. Therefore $[[0,\mu_2,\mu_3]] \notin R$, so this cube is in the region associated to some $Q \in V'$, where $Q \neq (b_{11},0,0)$. Since $[[0,\mu_2,\mu_3]]$ is labeled with $b$, we can apply the same argument as in the earlier cases. \qedhere
\end{itemize}
\end{proof}





\begin{lemma}\label{lemma_catenary_lambda1_secondary}
    Suppose that $S$ is secondary.  If $(-\lambda_1,\lambda_2,\lambda_3) \in V$ satisfies
    $$ \lambda_0d_0 + \lambda_1d_1 = \lambda_2d_2 + \lambda_3d_3$$
    with $\lambda_i \in \mathbb{Z}_{>0}$, then $\lambda_0d_0 + \lambda_1d_1 \in \textup{Betti}(S)$ and $\varphi^{-1}(\lambda_0d_0 + \lambda_1d_1) = \mathcal{F}_1 \cup \mathcal{F}_2$, where 
    \begin{align*}
        &\mathcal{F}_1 =  \{(\lambda_0 + \lambda a_{00},\lambda_1 - \lambda a_{11},0,0) \mid \lambda \in \mathbb{Z}, \lambda_0 + \lambda a_{00} \geq 0,\lambda_1 - \lambda a_{11} \geq 0\}, \\
        &\mathcal{F}_2 =  \{(0,0,\lambda_2 + \lambda a_{22},\lambda_3 - \lambda a_{33}) \mid \lambda \in \mathbb{Z}, \lambda_2 + \lambda a_{22} \geq 0,\lambda_3 - \lambda a_{33} \geq 0\}.
    \end{align*}
  
\end{lemma}

\begin{proof}
    Clearly $\mathcal{F}_1 \cup \mathcal{F}_2 \subseteq \varphi^{-1}(\lambda_0d_0 + \lambda_1d_1)$.
    Consider any $(\mu_0,\mu_1,\mu_2,\mu_3) \in \varphi^{-1}(\lambda_0d_0 + \lambda_1d_1)$.

    \begin{itemize}
    \item 

    Suppose $\mu_0 > 0$. If $\mu_2 > 0$, then by Lemma \ref{lemma_procedure_1}, no cube labeled with an element of the Ap{\'e}ry set is in the region associated to $(-\mu_1,\lambda_2-\mu_2,\lambda_3-\mu_3)$. This contradicts the fact $[[0,\lambda_2-1, \lambda_3]] \in T$, which holds since $(-\lambda_1,\lambda_2,\lambda_3) \in V$. If $\mu_3 > 0$, we get an analogous contradiction, hence $\mu_2=\mu_3 = 0$. If $(\mu_0,\mu_1,0,0) \notin \mathcal{F}_1$, then there exists $\lambda \in \mathbb{Z}$ such that $\lambda_0 + (\lambda+1) a_{00} > \mu_0  > \lambda_0 + \lambda a_{00}$. Now, we have
    \begin{equation}
        \lambda_0d_0 + \lambda_1d_1 = (\mu_0 - \lambda a_{00})d_0 + (\mu_1 + \lambda a_{11})d_1 \implies (\mu_0  - \lambda_0- \lambda a_{00})d_0 = (\lambda_1 - \mu_1 - \lambda a_{11})d_1
    \end{equation}
    which gives a contradiction, since $a_{00} > (\mu_0 - \lambda_0 - \lambda a_{00}) > 0$.

    \item
    Suppose $\mu_0 = 0$ and $\mu_1 > 0$. If $\mu_2 \geq \lambda_2$, then we get
    \begin{equation}
         (\lambda_3-\mu_3)d_3=\mu_1d_1 + (\mu_2 - \lambda_2)d_2
    \end{equation}
    where all coefficients above are nonnegative. Moreover, we have $(\lambda_3-\mu_3) > 0$, as $\mu_1 > 0$. Hence $(\lambda_3-\mu_3) \geq a_{33}$, so there exists $\lambda \in \mathbb{Z}_{>0}$ such that $a_{33} >(\lambda_3-\mu_3 - \lambda a_{33}) \geq 0$. We have 
    \begin{equation}
         (\lambda_3-\mu_3-\lambda a_{33})d_3=\mu_1d_1 + (\mu_2 - \lambda_2-\lambda a_{22})d_2.
    \end{equation}
    If $(\mu_2 - \lambda_2-\lambda a_{22}) \geq 0$, then as $\mu_1 > 0$ we have $(\lambda_3-\mu_3-\lambda a_{33}) > 0$, which gives $(\lambda_3-\mu_3-\lambda a_{33}) \geq a_{33}$ --- contradiction. Hence $(\mu_2 - \lambda_2-\lambda a_{22}) < 0$, which gives $\mu_1 \geq a_{11}$. Then, repeating the approach of the first bullet-point to the factorization $(a_{00},\mu_1 - a_{11},\mu_2,\mu_3) \in \varphi^{-1}(\lambda_0d_0 + \lambda_1d_1)$, yields $\mu_2 = \mu_3 = 0$ and $(a_{00},\mu_1 - a_{11},0,0) \in \mathcal{F}_1$, so $(\mu_0,\mu_1,\mu_2,\mu_3) =(0,\mu_1,0,0) \in \mathcal{F}_1$ also.
    
    If $\mu_3 \geq \lambda_3$, we proceed analogously; therefore, suppose that $\mu_2 < \lambda_2$ and $\mu_3 < \lambda_3$. We have 
    \begin{equation}\label{lemma_catenary_lambda1latter_secondary_eq1}
        \lambda_0d_0 = (\mu_1-\lambda_1)d_1 + \mu_2d_2 + \mu_3d_3. 
    \end{equation} 
     By Lemma \ref{lemma_procedure_1}, no cube labeled with an element of the Ap{\'e}ry set is in the region associated to $(\mu_1 - \lambda_1,\mu_2,\mu_3)$. If $(\mu_1 - \lambda_1) \leq 0$, this gives a contradiction with the fact $[[0,\lambda_2-1,\lambda_3-1]] \in T$, which holds since $(-\lambda_1,\lambda_2,\lambda_3) \in V$. Hence $(\mu_1 - \lambda_1) > 0$, which gives $\lambda_0 \geq a_{00}$. If $\lambda_0 = a_{00}$, then $\mu_2 = \mu_3 = 0$ and $\mu_1 = \lambda_1 + a_{11}$, since $\varphi^{-1}(a_{00}d_0) = \{(a_{00},0,0,0),(0,a_{11},0,0)\}$. This gives $(\mu_0,\mu_1,\mu_2,\mu_3) \in \mathcal{F}_1$, so suppose $\lambda_0 > a_{00}$. We have
    \begin{equation}
        (\lambda_0-a_{00})d_0 = (\mu_1-\lambda_1-a_{11})d_1 + \mu_2d_2 + \mu_3d_3. 
    \end{equation}
    Now, we can repeat the same argument used for \eqref{lemma_catenary_lambda1latter_secondary_eq1} to eventually obtain a contradiction or $(\mu_0,\mu_1,\mu_2,\mu_3) \in \mathcal{F}_1$. 
    
    \item 
    Suppose $\mu_0 =\mu_1 = 0$. If $(0,0,\mu_2,\mu_3) \notin \mathcal{F}_2$, then there exists $\lambda \in \mathbb{Z}$ such that $ \lambda_2  + (\lambda+1) a_{22} > \mu_2 > \lambda_2  + \lambda a_{22}$. Now, we have
    \begin{equation}
        \lambda_2d_2 + \lambda_3d_3 = (\mu_2 - \lambda a_{22})d_2 + (\mu_3 + \lambda a_{33})d_3 \implies (\mu_2 - \lambda_2 - \lambda a_{22} )d_2 = (\lambda_3 - \mu_3 - \lambda a_{33})d_3
    \end{equation}
    which gives a contradiction, since $ a_{22} > (\mu_2 - \lambda_2  - \lambda a_{22}) > 0$. \qedhere
     \end{itemize}
\end{proof}


\begin{lemma}\label{lemma_catenary_lambda23_secondary}

    Suppose that $S$ is secondary.  Consider $(\lambda_1,-\lambda_2,\lambda_3) \in V$ satisfying
    $$ \lambda_0d_0 + \lambda_2d_2 = \lambda_1d_1 + \lambda_3d_3$$
    with $\lambda_i \in \mathbb{Z}_{>0}$. If $\lambda_0 \geq a_{00}$ or $\lambda_2 \geq a_{22}$, then $\lambda_0d_0 + \lambda_2d_2 \notin \textup{Betti}(S)$.
    Otherwise, if $\lambda_0 < a_{00}$ and $\lambda_2 < a_{22}$, then $\lambda_0d_0 + \lambda_2d_2 \in \textup{Betti}(S)$ and $\varphi^{-1}(\lambda_0d_0 + \lambda_2d_2) = \{(\lambda_0,0,\lambda_2,0),(0,\lambda_1,0,\lambda_3)\}$. An analogous statement holds for $(\lambda_1,\lambda_2,-\lambda_3) \in V$.
\end{lemma}

\begin{proof}
    If $\lambda_0 \geq a_{00}$, then $(\lambda_0,0,\lambda_2,0) \, \mathcal{R} \, (\lambda_0-a_{00},a_{11},\lambda_2,0) \, \mathcal{R} \, (0,\lambda_1,0,\lambda_3)$, which implies that $\lambda_0d_0 + \lambda_2d_2$ has only one $\mathcal{R}$-class. If $\lambda_2 \geq a_{22}$ we can proceed analogously. Thus the first part of the lemma follows.

    Suppose that $\lambda_0 < a_{00}$ and $\lambda_2 < a_{22}$. Consider any $(\mu_0,\mu_1,\mu_2,\mu_3) \in \varphi^{-1}(\lambda_0d_0 + \lambda_2d_2)$.

    \begin{itemize}
        \item Suppose $\mu_0>0$. If $\mu_1 > 0$, then by Lemma \ref{lemma_procedure_1}, no cube labeled with an element of the Ap{\'e}ry set is in the region associated to $(\lambda_1-\mu_1,-\mu_2,\lambda_3-\mu_3)$. This contradicts the fact $[[\lambda_1-1,0, \lambda_3]] \in T$, which holds since $(\lambda_1,-\lambda_2,\lambda_3) \in V$. If $\mu_3 > 0$, we get an analogous contradiction, hence $\mu_1=\mu_3 = 0$. This gives
        \begin{equation}
            (\lambda_0- \mu_0) d_0 =  (\mu_2 - \lambda_2)d_2.
        \end{equation}
        If $(\lambda_0- \mu_0) = 0$, then $(\mu_0,\mu_1,\mu_2,\mu_3) = (\lambda_0,0,\lambda_2,0)$. If $(\lambda_0- \mu_0) > 0$, then $\lambda_0 \geq (\lambda_0- \mu_0) \geq a_{00}$ --- contradiction. If $(\lambda_0- \mu_0) < 0$, then $\lambda_2 \geq (\lambda_2- \mu_2) \geq a_{22}$ --- contradiction.


        \item
    Suppose $\mu_0 = 0$ and $\mu_2 > 0$. If $\mu_1 \geq \lambda_1$, then
    \begin{equation}
         (\lambda_3-\mu_3)d_3 = (\mu_1-\lambda_1)d_1 + \mu_2d_2
    \end{equation}
    where all coefficients above are nonnegative. If $(\lambda_3-\mu_3) = 0$, then $(\mu_0,\mu_1,\mu_2,\mu_3) = (0,\lambda_1,0,\lambda_3)$. Thus, suppose $(\lambda_3-\mu_3) > 0$, which implies $\lambda_3 \geq (\lambda_3-\mu_3) \geq a_{33}$. Now, as $\lambda_2 < a_{22}$, we have
    \begin{equation}\label{lemma_catenary_lambda23_secondary_eq2}
        \lambda_0d_0 = \lambda_1d_1 + (a_{22}-\lambda_2)d_2 + (\lambda_3 - a_{33})d_3
    \end{equation}
    where all coefficients above are nonnegative, which contradicts $\lambda_0 < a_{00}$. 
    
    If $\mu_3 \geq \lambda_3$, then
    \begin{equation}
         (\lambda_1-\mu_1)d_1 = (\mu_3-\lambda_3)d_3 + \mu_2d_2
    \end{equation}
    where all coefficients above are nonnegative. If $(\lambda_1-\mu_1) = 0$, then $(\mu_0,\mu_1,\mu_2,\mu_3) = (0,\lambda_1,0,\lambda_3)$. Thus, suppose $(\lambda_1-\mu_1) > 0$, which implies $\lambda_1 \geq (\lambda_1-\mu_1) \geq a_{11}$. Now, as $\lambda_0 < a_{00}$, we have
    \begin{equation}\label{lemma_catenary_lambda23_secondary_eq1}
        \lambda_2d_2 = (a_{00}-\lambda_0)d_0 + (\lambda_1 - a_{11})d_1 + \lambda_3d_3
    \end{equation}
    where all coefficients above are nonnegative, which contradicts $\lambda_2 < a_{22}$.

    Hence $\mu_1 < \lambda_1$ and $\mu_3 < \lambda_3$. We have
    \begin{equation}
         \lambda_0d_0 = \mu_1d_1 + (\mu_2 - \lambda_2)d_2 + \mu_3d_3.
    \end{equation}
     By Lemma \ref{lemma_procedure_1}, no cube labeled with an element of the Ap{\'e}ry set is in the region associated to $(\mu_1,\mu_2-\lambda_2,\mu_3)$. If $(\mu_2 - \lambda_2) \leq 0$, this gives a contradiction with the fact $[[\lambda_1-1,0,\lambda_3-1]] \in T$, which holds since $(\lambda_1,-\lambda_2,\lambda_3) \in V$. Hence $(\mu_2 - \lambda_2) > 0$, which gives $\lambda_0 \geq a_{00}$ --- contradiction.

    \item 
    Finally, suppose $\mu_0 = \mu_2 = 0$. We have
    \begin{equation}
        (\lambda_1 - \mu_1)d_1 = (\mu_3 - \lambda_3)d_3.
    \end{equation}
    If $(\lambda_1 - \mu_1) = 0$, then $(0,\mu_1,0,\mu_3) = (0,\lambda_1,0,\lambda_3)$. 
    If $(\lambda_1 - \mu_1) > 0$, then $\lambda_1 \geq (\lambda_1 - \mu_1) \geq a_{11}$ and we obtain \eqref{lemma_catenary_lambda23_secondary_eq1} with the same contradiction.
    Hence $(\lambda_1 - \mu_1) < 0$, which implies $\lambda_3 \geq (\lambda_3 - \mu_3) \geq a_{33} $. In this case we obtain \eqref{lemma_catenary_lambda23_secondary_eq2} leading to the same contradiction.
    \end{itemize}



The same argument applies to $(\lambda_1,\lambda_2,-\lambda_3) \in V$. 
\end{proof}

From Lemmas \ref{lemma_catenary_lambda1_secondary} and \ref{lemma_catenary_lambda23_secondary}, we know which elements of $U$ corresponding to $(2,2)$-type points are Betti elements. Moreover, we get that every Betti element that corresponds to a $(2,2)$-type point has exactly two $\mathcal{R}$-classes. We also know that $a_{ii}d_i \in \textup{Betti}(S)$ for $i = 0,1,2,3$ by Lemma \ref{lemma_catenary_minrelations}. However, we do not know whether $b_{22}d_2$ and $b_{33}d_3$ are Betti elements. This gap will not affect the results of Section \ref{section_type}, so we leave it unresolved.



\begin{lemma}\label{lemma_catenaryT_secondary_case1}
    Suppose that $S$ is secondary. No element $(\lambda_1,-\lambda_2,\lambda_3) \in V$ satisfies
    $$ \lambda_0d_0 + \lambda_2d_2 = \lambda_1d_1 + \lambda_3d_3 \notin \textup{Betti}(S)$$
    with $\lambda_i \in \mathbb{Z}_{>0}$.
\end{lemma}
\begin{proof}
By Lemma \ref{lemma_catenary_lambda23_secondary}, we have $\lambda_0 \geq a_{00}$ or $\lambda_2 \geq a_{22}$. Since $(\lambda_1,-\lambda_2,\lambda_3) \in V$ and $T = R \cap \{z \leq a_{33}\}$, we have $\lambda_1 < b_{11} = a_{11}$ and $\lambda_3 < a_{33}$. 

If $\lambda_0\geq a_{00}$, then 
\begin{equation}
     (\lambda_0-a_{00})d_0 + \lambda_2d_2 + (a_{11}-\lambda_1)d_1 =  \lambda_3d_3 
\end{equation}
where all above coefficients are nonnegative. Since $(a_{11}-\lambda_1) > 0$, we get $\lambda_3 \geq a_{33}$, which is a contradiction.

If $\lambda_2\geq a_{22}$, then
\begin{equation}
     \lambda_0d_0 + (\lambda_2-a_{22})d_2 + (a_{33}-\lambda_3)d_3 =  \lambda_1d_1 
\end{equation}
where all above coefficients are nonnegative. Since $(a_{33}-\lambda_3) > 0$, we get $\lambda_1 \geq a_{11}$, which is a contradiction.
\end{proof}

\begin{lemma}\label{lemma_catenaryT_secondary_case2}
    Suppose that $S$ is secondary. If $(\lambda_1,\lambda_2,-\lambda_3) \in V$ satisfies
    $$ \lambda_0d_0 + \lambda_3d_3 = \lambda_1d_1 + \lambda_2d_2 \notin \textup{Betti}(S)$$
    with $\lambda_i \in \mathbb{Z}_{>0}$, then $\lambda_0 = a_{00}$ and $\lambda_3 < a_{33}$.
\end{lemma}

\begin{proof}
By Lemma \ref{lemma_catenary_lambda23_secondary}, we have $\lambda_0 \geq a_{00}$ or $\lambda_3 \geq a_{33}$.  Since $(\lambda_1,-\lambda_2,\lambda_3) \in V$ and $T = R \cap \{z \leq a_{33}\}$, we have $\lambda_1 < b_{11} = a_{11}$.

If $\lambda_3\geq a_{33}$, then
\begin{equation}
     \lambda_0d_0   + (\lambda_3-a_{33})d_3=  \lambda_1d_1 + (\lambda_2-a_{22})d_2 
\end{equation}
where $(\lambda_3-a_{33}) \geq 0$.
If $(\lambda_2-a_{22}) < 0$, then $\lambda_1 \geq a_{11}$, which is a contradiction. Hence $(\lambda_2-a_{22}) \geq 0$. Since $\lambda_0 > 0$, no cube labeled with an element of the Ap{\'e}ry set is in the region associated to $(\lambda_1,\lambda_2-a_{22},-\lambda_3+a_{33})$ by Lemma \ref{lemma_procedure_1}. This contradicts the fact $[[\lambda_1,\lambda_2-1,0]] \in T$, which holds since $(\lambda_1,\lambda_2,-\lambda_3) \in V$. Hence $\lambda_3 < a_{33}$.

If $\lambda_0 > a_{00}$, then
\begin{equation}
     (\lambda_0-a_{00})d_0 + \lambda_3d_3 + (a_{11}-\lambda_1)d_1 =  \lambda_2d_2
\end{equation}
where all above coefficients are positive. This gives $\lambda_2 \geq b_{22}$, which contradicts $(\lambda_1,\lambda_2,-\lambda_3) \in V$. Hence $\lambda_0 = a_{00}$.
\end{proof}







\section{$t(S)$ and $\eta(S)$ }\label{section_type}


This section is devoted to proving Theorem \ref{theorem_type_main}. We begin with a series of lemmas relating the type to the set $V$, culminating in Theorem \ref{theorem_type1}. We then turn to results linking $V$ to $\eta(S)$, concluding with Corollary \ref{corollary_type_main}, from which Theorem \ref{theorem_type_main} follows.

Let $S = \langle d_0,d_1,d_2,d_3 \rangle$ be a numerical semigroup with embedding dimension four that is ordered. Let $T$, $R$, $V$, and $U$ be as in the previous sections.


\subsection{$t(S)$ and $V$}

A finite and nonempty subset $F$ of the initial collection of cubes is called a \emph{cube figure} if for every $(a,b,c) \in \mathbb{N}^3$ with $[[a,b,c]] \notin F$, no cube in the region associated to $(a,b,c)$ is in $F$. For a cube figure $F$, a cube $[[a,b,c]] \in F$ is called a \emph{corner} of $F$ if $[[a+1,b,c]],[[a,b+1,c]],[[a,b,c+1]] \notin F$. 

Recall from \eqref{eq_pseudo_frob} that
\begin{equation}
     PF(S) = \{s - d_0 \mid s \in \textup{Ap}(S,d_0), s + d_i \notin \textup{Ap}(S,d_0) \text{ for } i = 1,2,3\}.
\end{equation}
Thus, since $R$ is the collection of all cubes labeled with an element of the Ap{\'e}ry set, $PF(S)$ is the set of all labels of corners of $R$, and $t(S)$ is the number of such distinct labels. 







\begin{lemma}\label{lemma_point0uniqueness}
    There exists at most one $(0,a_{01},a_{02},a_{03}) \in \varphi^{-1}(a_{00}d_0)$ that satisfies $(a_{01},a_{02},a_{03}) \in V \setminus \{(b_{11},0,0),(0,b_{22},0),(0,0,b_{33})\}$.
\end{lemma}
\begin{proof}
If $S$ is secondary, then $\varphi^{-1}(a_{00}d_0) = \{(a_{00},0,0,0), (0,a_{11},0,0)\}$ (and by definition $a_{11}=b_{11}$), so no such $(0,a_{01},a_{02},a_{03}) \in \varphi^{-1}(a_{00}d_0)$ exists in this case. 

Suppose that $S$ is primary. Suppose to the contrary, that $(0,a_{01},a_{02},a_{03})$ and $(0,a_{01}',a_{02}',a_{03}') \in \varphi^{-1}(a_{00}d_0)$ are two such distinct factorizations. Since $T \subseteq \{x \leq a_{11}, y \leq a_{22}, z \leq a_{33}\}$, it follows that $a_{0i}, a_{0i}' \leq a_{ii}$. If $(a_{01},a_{02},a_{03}) = (0,0,a_{33})$, then $a_{33}=b_{33}$ (as $a_{00}>0$), which contradicts $(a_{01},a_{02},a_{03}) \neq (0,0,b_{33})$. If  $(a_{01},a_{02},a_{03}) \in \{(a_{11},0,0) ,(0,a_{22},0) \}$, we get an analogous contradiction, therefore $a_{0i}, a_{0i}' < a_{ii}$.  We have
\begin{equation}
    (a_{01}-a_{01}')d_1 +  (a_{02}-a_{02}')d_2 + (a_{03}-a_{03}')d_3 = 0.
\end{equation}
Since $(0,a_{01},a_{02},a_{03}) \neq (0,a_{01}',a_{02}',a_{03}')$, not all coefficients above are zero.
Without loss of generality, suppose $(a_{01}-a_{01}') >0$, $(a_{02}-a_{02}') \leq 0$, and $(a_{03}-a_{03}') \leq 0$. This implies $a_{01} \geq (a_{01}-a_{01}') \geq a_{11}$, which is a contradiction.
\end{proof}

Let $V^*$ denote the set of all $(2,2)$-type points of $V$. By the definition of $V$ and Lemma \ref{lemma_point0uniqueness}, we get
\begin{equation}\label{type_assumption_on_V}
    V = V^* \sqcup \{(b_{11},-b_{12},-b_{13}), (-b_{21},b_{22},-b_{23}), (-a_{31},-a_{32},a_{33}), (a_{01},a_{02},a_{03})\}
\end{equation}
where $a_{ij},b_{ij} \in \mathbb{N}$ and $(0,a_{01},a_{02},a_{03}) \in \varphi^{-1}(a_{00}d_0)$. Such a point $(a_{01},a_{02},a_{03}) \in V$ exists, since $a_{00}d_0 \in \textup{Betti}(S) \subseteq U \cup \{b_{33}d_3\}$ by Theorems \ref{theorembettifind_primary} and \ref{theorembettifind_secondary} --- perhaps $(a_{01},a_{02},a_{03}) \in \{(b_{11},-b_{12},-b_{13})$, $(-b_{21},b_{22},-b_{23})$, $(-a_{31},-a_{32},a_{33}) \}$. 

The identity \eqref{type_assumption_on_V} yields the following lemma.



\begin{lemma}\label{lemma_point0connection}
    If $(\lambda_1,\lambda_2,\lambda_3) \in V$ with $\lambda_i \geq 0$ for $i=1,2,3$ and $\lambda_i > 0$ for at least two distinct indices $i \in \{1,2,3\}$, then $(0,\lambda_1,\lambda_2,\lambda_3) \in \varphi^{-1}(a_{00}d_0)$.
\end{lemma}









\begin{lemma}\label{lemma_cornerlabelRT}
    Every integer that labels a corner of $R$ also labels a corner of $T$.
\end{lemma}

\begin{proof}
Take $(a_{30},a_{31},a_{32},0) \in \varphi^{-1}(a_{33}d_3)$. If $a_{30}>0$, then $a_{33}=b_{33}$ and $R = T$, which yields the claim. Thus, suppose $a_{30} =0$.

Consider a corner $[[a,b,c]]$ of $R$ that is in the region $R \cap \{z \geq a_{33}\}$ (that is $c \geq a_{33}$). There exists $\lambda \in \mathbb{N}$ such that $a_{33} > c- \lambda a_{33} \geq 0$, for which we have $C = [[a + \lambda a_{31},b + \lambda a_{32},c- \lambda a_{33}]] \in  R \cap \{z \leq a_{33}\} = T$, since $C$ has the same label as $[[a,b,c]]$. 
We claim that $C$ is a corner of $T$. If not, then there exists $[[a',b',c']] \in T $ with $C <  [[a',b',c']]$. However, this implies $[[a,b,c]] < [[a' -\lambda a_{31},b'-\lambda a_{32},c' + \lambda a_{33} ]] \in R$, which contradicts the fact that $[[a,b,c]]$ is a corner of $R$.
\end{proof}




Let $t'$ denote the number of corners of $T$ and let $\mathcal{T}$ be the set of the corners of $T$ that are not corners of $R$. Since $T$ is an $L$-shape, the corners of $T$ have distinct labels. This fact, combined with Lemma \ref{lemma_cornerlabelRT}, yields $t' = | \mathcal{T}| + t(S)$.




For the semigroups $\langle 196,225,256,289 \rangle$ and $\langle 145,203,74,111 \rangle$ from Example \ref{example_T}, shown in Figures \ref{fig_cornersT_primary} and \ref{fig_cornersT_secondary}, respectively, the cubes shown in blue are all the corners of $T$ that are also corners of $R$, and the cubes shown in magenta are all the cubes in $\mathcal{T}$.


\begin{figure}[h]
    \centering
    \begin{minipage}[t]{0.48\textwidth}
        \centering
        \includegraphics[width=0.7\textwidth]{tintar_cornersT_primary.pdf}
        \caption{}
        \label{fig_cornersT_primary}
    \end{minipage}
    \begin{minipage}[t]{0.48\textwidth}
        \centering
        \includegraphics[width=0.7\textwidth]{tintar_cornersT_secondary.pdf}
        \caption{}
        \label{fig_cornersT_secondary}
    \end{minipage}
\end{figure}



\begin{lemma}\label{lemma_type-corners_primary}
    If $S$ is primary, then $2t(S) \geq t' \geq t(S)$.
\end{lemma}
\begin{proof}
If $a_{33} = b_{33}$, then $t' = t(S)$, so suppose that $a_{33} < b_{33}$. Since $t' = | \mathcal{T}| + t(S)$, it follows immediately that $t' \geq  t(S)$, so we only need to prove $t' \leq 2t(S) $, which is equivalent to $|\mathcal{T}| \leq t(S)$.

Consider any $[[\lambda_1,\lambda_2,\lambda_3]] \in \mathcal{T}$. Since $[[\lambda_1,\lambda_2,\lambda_3]]$ is not a corner of $R$ and $T = R \cap \{z \leq a_{33}\} $, we have $\lambda_3 = a_{33}-1$ and $[[\lambda_1,\lambda_2,a_{33}]] \in R$. There exists $k \in \mathbb{N}$ such that $[[\lambda_1,\lambda_2,a_{33}+k]] \in R$ and $[[\lambda_1,\lambda_2,a_{33}+k+1]] \notin R$. Note that $[[\lambda_1,\lambda_2,a_{33}+k]]$ is a corner of $R$, since $T = R \cap \{z \leq a_{33}\}$ and $[[\lambda_1,\lambda_2,a_{33}-1]]$ is a corner of $T$. 

If  $| \mathcal{T}| > t(S)$, then there exist two distinct cubes $[[\lambda_1,\lambda_2,a_{33}-1]], [[\lambda_1',\lambda_2',a_{33}-1]] \in \mathcal{T}$ such that, for some $k, k' \in \mathbb{N}$, $[[\lambda_1,\lambda_2,a_{33}+k]]$ and $[[\lambda_1',\lambda_2',a_{33}+k']]$ are corners of $R$ with the same label. Without loss of generality, suppose $k \geq k'$, $\lambda_1 > \lambda_1'$ and $\lambda_2 < \lambda_2'$ (if $\lambda_1 \geq \lambda_1'$ and $\lambda_2 \geq \lambda_2'$, we would get a contradiction with $[[\lambda_1,\lambda_2,a_{33}-1]]$ and $[[\lambda_1',\lambda_2',a_{33}-1]]$ being distinct corners of $T$). This gives
\begin{equation}
    (\lambda_1 - \lambda_1')d_1 + (k-k')d_3 = (\lambda_2'-\lambda_2)d_2
\end{equation}
which implies $\lambda_2' \geq (\lambda_2'-\lambda_2) \geq a_{22} = b_{22}$, contradicting $[[\lambda_1',\lambda_2',a_{33}-1]] \in T$. Thus $|\mathcal{T}| \leq t(S)$.    
\end{proof}

\begin{lemma}\label{lemma_type-corners_secondary}
    If $S$ is secondary, then $t(S)+1 \geq t' \geq  t(S)$.
\end{lemma}

\begin{proof}
    As before, $t' = | \mathcal{T}| + t(S)$ gives $t' \geq  t(S)$, so it remains to show $t' \leq t(S)+1 $, which is equivalent to $|\mathcal{T}| \leq 1$.

    Consider any $[[\lambda_1,\lambda_2,\lambda_3]] \in \mathcal{T}$. Since $[[\lambda_1,\lambda_2,\lambda_3]]$ is not a corner of $R$ and $T = R \cap \{z \leq a_{33}\} $, we have $\lambda_3 = a_{33}-1$ and $[[\lambda_1,\lambda_2,a_{33}]] \in R$.
    
    If $|\mathcal{T}| \geq 2$, then there exist $[[\lambda_1,\lambda_2,a_{33}-1]],  [[\lambda_1',\lambda_2',a_{33}-1]] \in \mathcal{T}$ with $\lambda_1 > \lambda_1'$ and $\lambda_2 < \lambda_2'$. Additionally, assume that no cube $[[\lambda_1'',\lambda_2'',a_{33}-1]] \in \mathcal{T}$ satisfies $\lambda_1 > \lambda_1'' > \lambda_1'$ and $\lambda_2 < \lambda_2'' < \lambda_2'$. This assumption is equivalent to $(\lambda_1'+1,\lambda_2+1,\lambda) \in V$ for some $\lambda \in \mathbb{Z}$. 
    
    If $\lambda \geq 0$, then by Lemma \ref{lemma_point0connection}, as $\lambda_1'+1,\lambda_2+1 > 0$, we get $(0,\lambda_1'+1,\lambda_2+1,\lambda) \in \varphi^{-1}(a_{00}d_0)$. As $S$ is secondary, we have $\varphi^{-1}(a_{00}d_0) = \{(a_{00},0,0,0),(0,a_{11},0,0)\}$, which is a contradiction. 
    
    Hence $\lambda < 0$, which implies that no cube in the region $\{x \geq \lambda_1'+1, y \geq  \lambda_2+1, z \geq 0\}$ is in $R$. We have $[[\lambda_1,\lambda_2,a_{33}]] \in R$, since $[[\lambda_1,\lambda_2,a_{33}-1]] \in \mathcal{T}$. By Theorem \ref{theoremprocedure}, this gives $[[\lambda_1,\lambda_2+a_{22},0]] \in R$, as $a_{22}d_2 = a_{33}d_3$. However, $[[\lambda_1,\lambda_2+a_{22},0]] \in \{x \geq \lambda_1'+1, y \geq  \lambda_2+1, z \geq 0\}$, as $\lambda_1 > \lambda_1'$, which is a contradiction. Thus $|\mathcal{T}| \leq 1$.   
\end{proof}












Let $F$ be a cube figure.
A cube $[[a,b,0]] \in F$ with $[[a+1,b,0]], [[a,b+1,0]] \notin F$ is called a \emph{stair cube}. Similarly, $[[a,0,c]] \in F$ with $[[a+1,0,c]], [[a,0,c+1]] \notin F$, and $[[0,b,c]] \in F$ with $[[0,b+1,c]], [[0,b,c+1]] \notin F$ are also \emph{stair cubes}.  

Let $\mathcal{S}$ denote the set of all stair cubes of $T$. For the primary semigroup $\langle 784, 841,900,961 \rangle$, Figure \ref{fig_staircubes} shows the corresponding figure $T$ with all stair cubes shown in blue.




\begin{figure}[h]
    \centering
    \includegraphics[width=0.5\textwidth]{tintar_staircubes.pdf}
    \caption{}
    \label{fig_staircubes}
\end{figure}


\begin{lemma}\label{lemma_type_suppthm_stair}
 If there exists $(0,a_{01},a_{02},a_{03}) \in \varphi^{-1}(a_{00}d_0)$ such that $(a_{01},a_{02},a_{03}) \in V$ and exactly one among $a_{01}, a_{02}, a_{03}$ is zero, then $|V^*|+4 \geq |\mathcal{S}| \geq |V^*|+1$. Otherwise $|V^*|+3 \geq |\mathcal{S}| \geq |V^*|$.
\end{lemma}

\begin{proof}
     Order the elements of $\mathcal{S}$ with $z=0$ according to increasing $x$-coordinate. Note that they are also ordered by decreasing $y$-coordinate. In this order, consider two consecutive stair cubes with $z=0$: $[[\lambda_1,\lambda_2,0]], [[\mu_1,\mu_2,0]]$ with $\lambda_1 > \mu_1$ and $\lambda_2 < \mu_2$. Since they are consecutive, there exists $\nu \in \mathbb{N}$ such that $(\mu_1+1,\lambda_2+1,-\nu) \in V$. 

     If no $(0,a_{01},a_{02},a_{03}) \in \varphi^{-1}(a_{00}d_0)$ satisfies $(a_{01},a_{02},a_{03}) \in V$ where exactly one among $a_{01}, a_{02}, a_{03}$ is zero, then $\nu > 0$, since otherwise $(\mu_1+1,\lambda_2+1,0) \in V$ would be such a point by Lemma \ref{lemma_point0connection}. This gives $(\mu_1+1,\lambda_2+1,-\nu) \in V^*$, which we assign to the stair cube $[[\lambda_1,\lambda_2,0]]$. Such a point $(\mu_1+1,\lambda_2+1,-\nu) \in V^*$ exists for all stair cubes in $T \cap \{z = 0\}$ except for the minimal one in our ordering. Moreover, every such point is unique (when $V$ is fixed), and it is clear that we assign exactly one stair cube of $T \cap \{z = 0\}$ to every point of $V^*$ in the region $\{x > 0, y > 0, z < 0\}$.   
     Therefore in the region $\{x > 0, y > 0, z < 0\}$ there is one fewer element of $V^*$ than there are stair cubes in $T \cap \{z = 0\}$. Applying the same approach to stair cubes of $T \cap \{x = 0\}$ and $T \cap \{y = 0\}$, we obtain $|V^*|+3 \geq |\mathcal{S}| \geq |V^*|$, since there are at most three stair cubes of $T$ that belong to more than one of the sets $T \cap \{x = 0\}$, $T \cap \{y = 0\}$, and $T \cap \{z = 0\}$; if such stair cubes exist, they must be among $[[b_{11}-1,0,0]], [[0,b_{22}-1,0]], [[0,0,a_{33}-1]]$.

     If there exists $(0,a_{01},a_{02},a_{03}) \in \varphi^{-1}(a_{00}d_0)$ such that $(a_{01},a_{02},a_{03}) \in V$ and exactly one among $a_{01}, a_{02}, a_{03}$ is zero, then this points in unique by Lemma \ref{lemma_point0uniqueness}. Without loss of generality, suppose $a_{03} = 0$. If $\nu = 0$, then necessarily $(a_{01},a_{02},a_{03}) = (\mu_1+1,\lambda_2+1,0) \in V$. Since $(a_{01},a_{02},a_{03})$ is unique, we get that in the region $\{x > 0, y > 0, z < 0\}$ there are two fewer elements of $V^*$ than there are stair cubes in $T \cap \{z = 0\}$. Again, since $(a_{01},a_{02},a_{03})$ is unique, applying the same approach to stair cubes of $T \cap \{x = 0\}$ and $T \cap \{y = 0\}$ gives the same result as in the previous paragraph. Therefore $|V^*|+4 \geq |\mathcal{S}| \geq |V^*|+1$.  
\end{proof}

\begin{lemma}\label{lemma_type_suppthm_upperbound}
    $  2t'+1 \geq |\mathcal{S}|$.
\end{lemma}

\begin{proof}
    Let $[[a,b,c]]$ be a corner of $T$. Consider how many cubes among $[[0,b,c]],[[a,0,c]],[[a,b,0]]$ are stair cubes of $T$. We claim that there is at most one corner of $T$ that has all three as stair cubes. Suppose to the contrary, that $[[a,b,c]]$ and $[[a',b',c']]$ are two such distinct corners of $T$. Without loss of generality, suppose $a > a'$, $b \geq b'$, $c < c'$. Now, since $[[a',b',0]]$ is a stair cube, we have $[[a'+1,b',0]] \notin T$. This implies $[[a,b,0]] \notin T$, which is a contradiction. 
    
    Therefore, assigning to each corner $[[a,b,c]]$ of $T$ the stair cubes among $[[0,b,c]],[[a,0,c]],$ $[[a,b,0]]$, every corner of $T$ receives at most two stair cubes, with the exception of at most one corner, which receives three. Moreover, each stair cube is assigned to exactly one corner of $T$, hence $ 2t'+1 \geq |\mathcal{S}|$.
\end{proof}



A corner $[[a,b,c]]$ of a cube figure $F$ is an \emph{odd corner}
if none of $[[0,b,c]]$, $[[a,0,c]]$, $[[a,b,0]]$ is a stair cube of $F$.


\begin{prop}\label{prop_type_suppthm_odd_corners}
    Let $F$ be a cube figure. If there is no $[[a,b,c]] \notin F$ such that $[[a-1,b,c]],[[a,b-1,c]],[[a,b,c-1]] \in F$, then $F$ has at most two odd corners.
\end{prop}

\begin{proof}
    Let $[[a,b,c]]$ be an odd corner of $F$. This implies: $[[a+1,b,0]] \in F$ or $[[a,b+1,0]] \in F$, $[[a+1,0,c]] \in F$ or $[[a,0,c+1]] \in F$, and $[[0,b+1,c]] \in F$ or $[[0,b,c+1]] \in F$. Suppose that $[[a+1,b,0]]$ and $[[a,b+1,0]]$ are both in $F$ --- we claim that this cannot be the case. 

    If $[[a+1,0,c]] \in F$, then there exists $b \geq \lambda > 0$ such that $[[a+1,\lambda-1,c]] \in F$ and $[[a+1,\lambda,c]] \notin F$. Then, there exists $c \geq \mu > 0$ such that $[[a+1,\lambda,\mu-1]] \in F$ and $[[a+1,\lambda,\mu]] \notin F$, as $[[a+1,\lambda,0]] \leq [[a+1,b,0]] \in F$. However, since $[[a+1,\lambda,\mu]] \notin F$ and $[[a,\lambda,\mu]], [[a+1,\lambda-1,\mu]], [[a+1,\lambda,\mu-1]] \in F$, we get a contradiction with the hypothesis. On Figure \ref{fig_type_theorem_1}, we can see an example of such a situation, where $P = (a+1,\lambda,\mu)$. 
    
    If $[[0,b+1,c]] \in F$, we get an analogous contradiction; thus, since $[[a,b,c]]$ is an odd corner and neither $[[a+1,0,c]]$ nor $[[0,b+1,c]]$ is in $F$, we have $[[a,0,c+1]], [[0,b,c+1]] \in F$. Now, there exist $a \geq \lambda > 0$, $b \geq \mu > 0$, such that $[[\lambda,\mu,c+1]] \notin F$ and $[[\lambda-1,\mu,c+1]], [[\lambda,\mu-1,c+1]] \in F$. Since $[[\lambda,\mu,c]] \in F$ also, we get a contradiction with the hypothesis. On Figure \ref{fig_type_theorem_2}, we can see an example of such a situation, where $P = (\lambda,\mu,c+1)$.


\begin{figure}[h]
    \centering
    \begin{minipage}[t]{0.48\textwidth}
        \centering
        \includegraphics[width=0.7\textwidth]{tintar_type_theorem_1.pdf}
        \caption{}
        \label{fig_type_theorem_1}
    \end{minipage}
    \begin{minipage}[t]{0.48\textwidth}
        \centering
        \includegraphics[width=0.7\textwidth]{tintar_type_theorem_2.pdf}
        \caption{}
        \label{fig_type_theorem_2}
    \end{minipage}
\end{figure}

    Similarly, we obtain that exactly one of $[[a+1,0,c]]$ and $[[a,0,c+1]]$ is in $F$, and that exactly one of $[[0,b+1,c]]$ and $[[0,b,c+1]]$ is in $F$. 

    Now, we claim that any two distinct odd corners $[[a,b,c]]$ and $[[a',b',c']]$ of $F$ must agree in at least one coordinate. Otherwise, without loss of generality, suppose that $a < a'$, $b < b'$, $c > c'$ (we cannot have $[[a,b,c]] < [[a',b',c']]$, as they are both corners). Then, clearly $[[a+1,b,0]], [[a,b+1,0]] \in F$, since $[[a+1,b,0]], [[a,b+1,0]] < [[a',b',c']] \in F$, which contradicts the claim proved at the beginning.

    
    Thus, without loss of generality, suppose that $a < a'$, $b = b'$, $c > c'$.  Suppose that $[[a'',b'',c'']]$ is a third distinct odd corner of $F$. If $b'' \neq b$, then from the previous paragraph
    we obtain $[[a'',b'',c'']] = [[a,b'',c']]$ or $[[a'',b'',c'']] = [[a',b'',c]]$ --- recall that $a < a'$ and $c > c'$, so we cannot have $a'' = a' = a$ or $c'' = c' = c$. 



    
    Suppose $[[a'',b'',c'']] = [[a',b'',c]]$. If $b'' \geq b' = b$, then $[[a',b'',c]] >[[a',b',c']]$, which contradicts the fact that $[[a',b',c']]$ is a corner of $F$. Hence $ b'' < b' = b$. Now, there exist $a' \geq \lambda > a$, $b \geq \mu > b''$, and $c \geq \nu > c'$, such that $[[\lambda,\mu,\nu]] \notin F$ and $[[\lambda-1,\mu,\nu]],[[\lambda,\mu-1,\nu]],[[\lambda,\mu,\nu-1]] \in F$, which contradicts the hypothesis. On Figure \ref{fig_type_theorem_3}, we can see an example of such a situation, where $P = (\lambda,\mu,\nu)$. 
    
    Suppose $[[a'',b'',c'']] = [[a,b'',c']]$. If $ b'' \leq b'=b $, then $[[a,b'',c']] < [[a,b,c]]$, which contradicts the fact that $ [[a,b'',c']]$ is a corner of $F$. Hence $b'' > b' =b$. Now, since $[[a+1,b,0]] \leq [[a',b',0]] \in F$ and $ [[a,b+1,0]] \leq [[a,b'',0]] = [[a'',b'',0]] \in F$, we have $[[a+1,b,0]], [[a,b+1,0]] \in F$, which contradicts the claim proved at the beginning. On Figure \ref{fig_type_theorem_4}, we can see an example of such a situation. 

\begin{figure}[h]
    \centering
    \begin{minipage}[t]{0.48\textwidth}
        \centering
        \includegraphics[width=0.7\textwidth]{tintar_type_theorem_3.pdf}
        \caption{}
        \label{fig_type_theorem_3}
    \end{minipage}
    \begin{minipage}[t]{0.48\textwidth}
        \centering
        \includegraphics[width=0.7\textwidth]{tintar_type_theorem_4.pdf}
        \caption{}
        \label{fig_type_theorem_4}
    \end{minipage}
\end{figure}


    
    Hence $b= b' = b'' $, and without loss of generality, suppose that
    $a < a' < a''$ and $c > c' > c''$. We have $[[a'+1,0,c']] \in F$ or $[[a',0,c'+1]] \in F$, as $[[a',b',c']]$ is an odd corner. 
    
    If $[[a',0,c'+1]] \in F$, then there exist $a' \geq \lambda > a$ and $b \geq \mu > 0$ such that $[[\lambda,\mu,c'+1]] \notin F$ and $[[\lambda-1,\mu,c'+1]], [[\lambda,\mu-1,c'+1]], [[\lambda,\mu,c']] \in F$, which contradicts the hypothesis. On Figure \ref{fig_type_theorem_5}, we can see an example of such a situation, where $P = (\lambda,\mu,c'+1)$.

     If $[[a'+1,0,c']] \in F$, then there exist $a'' \geq \lambda > a'$ and $b \geq \mu > 0$ such that $[[\lambda,\mu,c''+1]] \notin F$ and $[[\lambda-1,\mu,c''+1]], [[\lambda,\mu-1,c''+1]], [[\lambda,\mu,c'']] \in F$, which contradicts the hypothesis. On Figure \ref{fig_type_theorem_6}, we can see an example of such a situation, where $P = (\lambda,\mu,c'+1)$.  \qedhere
     


\begin{figure}[h]
    \centering
    \begin{minipage}[t]{0.48\textwidth}
        \centering
        \includegraphics[width=0.7\textwidth]{tintar_type_theorem_5.pdf}
        \caption{}
        \label{fig_type_theorem_5}
    \end{minipage}
    \begin{minipage}[t]{0.48\textwidth}
        \centering
        \includegraphics[width=0.7\textwidth]{tintar_type_theorem_6.pdf}
        \caption{}
        \label{fig_type_theorem_6}
    \end{minipage}
\end{figure}
\end{proof}


\begin{lemma}\label{lemma_type_suppthm_main_lower}
 If there exists $(0,a_{01},a_{02},a_{03}) \in \varphi^{-1}(a_{00}d_0)$ such that $(a_{01},a_{02},a_{03}) \in V$ and $a_{01} a_{02} a_{03}$ $ > 0$, then $|\mathcal{S}| \geq t'-8$. Otherwise $|\mathcal{S}| \geq t'-2$.
\end{lemma}
    

\begin{proof}

    If no $(0,a_{01},a_{02},a_{03}) \in \varphi^{-1}(a_{00}d_0)$ has $(a_{01},a_{02},a_{03}) \in V$ and $a_{01}a_{02}a_{03} > 0$, then $T$ satisfies the assumptions of Proposition \ref{prop_type_suppthm_odd_corners}, hence it has at most two odd corners. Then each non-odd corner $[[a,b,c]]$ of $T$ has at least one associated stair cube, chosen from $[[0,b,c]], [[a,0,c]],[[a,b,0]]$. Every stair cube is associated to at most one corner of $T$, hence $|\mathcal{S}| \geq t' -2$.
    
    
    Suppose that there exists $(0,a_{01},a_{02},a_{03}) \in \varphi^{-1}(a_{00}d_0)$ with $(a_{01},a_{02},a_{03}) \in V$ and $a_{01}a_{02}a_{03} > 0$. By Lemmas \ref{lemma_point0uniqueness} and \ref{lemma_point0connection} we know that this point is unique. Let $T'$ be the figure obtained by removing from the initial collection of cubes all the regions associated to each of $V \setminus \{ (a_{01},a_{02},a_{03})\}$. Note that the sets of stair cubes of $T$ and $T'$ are equal. Moreover, $T'$ satisfies the assumptions of Proposition \ref{prop_type_suppthm_odd_corners}, hence it has at most two odd corners.

    

    Now, deleting from $T'$ the region associated to $(a_{01},a_{02},a_{03})$ will yield $T$. We claim that this operation can add at most six new odd corners. Note that any new odd corner could not have been a corner of $T'$ --- by the definition, a corner that is not odd in $T'$ cannot become an odd corner merely because cubes were deleted. Thus every new odd corner must touch the region associated to $(a_{01},a_{02},a_{03})$ with one of its faces.

    Assume, for contradiction, that seven new odd corners appear. This means that three of them touch the same side of the region associated to $(a_{01},a_{02},a_{03})$. Without loss of generality, suppose that $[[x_1,y_1,a_{03}-1]], [[x_2,y_2,a_{03}-1]], [[x_3,y_3,a_{03}-1]]$ are three distinct odd corners, where $x_1 > x_2 > x_3 \geq a_{01}$ and $y_3 > y_2 > y_1 \geq a_{02}$. We have $[[x_2+1,y_2,0]] \in T$ or $[[x_2,y_2+1,0]] \in T$. If $[[x_2+1,y_2,0]] \in T$, then there exists $x_1 \geq \lambda > x_2$, $y_2 \geq  \mu > y_1$, and $a_{03}-1 \geq \nu > 0$, such that $[[\lambda,\mu,\nu]] \notin F$ and $[[\lambda-1,\mu,\nu]],[[\lambda,\mu-1,\nu]],[[\lambda,\mu,\nu-1]] \in F$. This implies $(\lambda,\mu,\nu) \in V$, which by Lemma \ref{lemma_point0connection}, contradicts Lemma \ref{lemma_point0uniqueness}. On Figure \ref{fig_type_theorem_7}, we can see an example of such a situation, where $P = (\lambda,\mu,\nu)$. If $[[x_2,y_2+1,0]] \in T$, we proceed analogously.

    \begin{figure}[h]
    \centering
    \includegraphics[width=0.5\textwidth]{tintar_type_theorem_7.pdf}
    \caption{}
    \label{fig_type_theorem_7}
    \end{figure}

    Thus, since $T'$ has at most two odd corners, $T$ can have at most eight odd corners.
    Each non-odd corner $[[a,b,c]]$ of $T$ has at least one associated stair cube, chosen from $[[0,b,c]], [[a,0,c]],[[a,b,0]]$. Every stair cube is associated to at most one corner of $T$, hence $|\mathcal{S}| \geq t' - 8$. 
\end{proof}


\begin{thm}\label{theorem_type1}
    If $S$ is primary, then $4 t(S) +1 \geq |V^*| \geq t(S)-11$.  If $S$ is secondary, then $2 t(S) +3 \geq |V^*| \geq t(S) - 5$.
\end{thm}


\begin{proof}
    Suppose that $S$ is primary. The upper bound for $|V^*|$ follows from the chain of inequalities:
    \begin{equation}
        |V^*| \leq |\mathcal{S}| \leq 2 t' + 1 \leq 4 t(S) +1
    \end{equation}
    which are the result of Lemmas \ref{lemma_type_suppthm_stair}, \ref{lemma_type_suppthm_upperbound}, and \ref{lemma_type-corners_primary}, respectively. To prove the lower bound for $|V^*|$ we need to split into two cases.

    Suppose that there exists $(0,a_{01},a_{02},a_{03}) \in \varphi^{-1}(a_{00}d_0)$ such that $(a_{01},a_{02},a_{03}) \in V$ and $a_{01} a_{02} a_{03} > 0$. Then, by Lemma \ref{lemma_point0uniqueness}, the point $(a_{01},a_{02},a_{03}) \in V$ is unique and there does not exist $(0,a_{01}',a_{02}',a_{03}') \in \varphi^{-1}(a_{00}d_0)$ such that $(a_{01}',a_{02}',a_{03}') \in V$ and exactly one among $a_{01}', a_{02}', a_{03}'$ is zero. Hence, we have the following chain of inequalities: 
    \begin{equation}
        |V^*| \geq |\mathcal{S}|-3 \geq t' - 11 \geq t(S) - 11
    \end{equation}
    which are the result of Lemmas \ref{lemma_type_suppthm_stair}, \ref{lemma_type_suppthm_main_lower}, and \ref{lemma_type-corners_primary}, respectively.

    Suppose that there does not exist $(0,a_{01},a_{02},a_{03}) \in \varphi^{-1}(a_{00}d_0)$ such that $(a_{01},a_{02},a_{03}) \in V$ and $a_{01} a_{02} a_{03} > 0$. Then, we have the following chain of inequalities: 
    \begin{equation}
        |V^*| \geq |\mathcal{S}|-4 \geq t' - 6 \geq t(S) - 6
    \end{equation}
    which are the result of Lemmas \ref{lemma_type_suppthm_stair}, \ref{lemma_type_suppthm_main_lower}, and \ref{lemma_type-corners_primary}, respectively. As $t(S) - 6 \geq t(S) - 11$, the first claim follows. 


    Suppose that $S$ is secondary. The upper bound for $|V^*|$ is due to the following chain of inequalities:
    \begin{equation}
        |V^*| \leq |\mathcal{S}| \leq 2 t' + 1 \leq 2 t(S) +3
    \end{equation}
    which are the result of Lemmas \ref{lemma_type_suppthm_stair}, \ref{lemma_type_suppthm_upperbound}, and \ref{lemma_type-corners_secondary}, respectively. Since $\varphi^{-1}(a_{00}d_0) = \{(a_{00},0,0,0)$, $(0,a_{11},0,0)\}$, there does not exist $(0,a_{01},a_{02},a_{03}) \in \varphi^{-1}(a_{00}d_0)$ such that $a_{01} a_{02} a_{03} > 0$, nor such that exactly one among $a_{01}, a_{02}, a_{03}$ is zero. Hence, we have the following chain of inequalities: 
    \begin{equation}
        |V^*| \geq |\mathcal{S}|-3 \geq t' -5 \geq t(S) - 5
    \end{equation} 
    which are the result of Lemmas \ref{lemma_type_suppthm_stair}, \ref{lemma_type_suppthm_main_lower}, and \ref{lemma_type-corners_secondary}, respectively.
\end{proof}














\subsection{$V$ and $\eta(S)$}
To relate $|V^*|$ with $\eta(S)$, we need to know how many elements of $V^*$ do not correspond to Betti elements. 

Let
$$ \mathcal{N} = \{ (x,y,z) \in V^* \mid  \max\{ x,0\} d_1 + \max\{y,0\}d_2 + \max\{z,0\}d_3 \notin \textup{Betti}(S)  \}.$$

\begin{lemma}\label{lemma_bettiT_primary}
     If $S$ is primary, then $|\mathcal{N}| \leq 3$.
\end{lemma}
\begin{proof} 
 

In this proof, as in the proof of Lemma \ref{lemmabettifind_primary}, it will be more convenient to drop the assumption that $S$ is ordered --- we only suppose that $a_{ii}=b_{ii}$ for at least two distinct $i \in \{1,2,3\}$ and that $T = R \cap \{x \leq a_{11},y\leq a_{22}, z \leq a_{33}\}$. 

Suppose $|\mathcal{N}| \geq 4$. Then, we can find two distinct points in $\mathcal{N}$ whose negative coordinates occur in the same position. Without loss of generality, let $(-\lambda_1,\lambda_2,\lambda_3)$ and $(-\lambda_1',\lambda_2',\lambda_3')$ be distinct points in $\mathcal{N}$ (where $\lambda_i>0$). By Lemma \ref{lemma_catenaryT_primary}, we have 
\begin{align}
    a_{00}d_0 + \lambda_1d_1 &=  \lambda_2d_2 + \lambda_3d_3,  \label{lemma_bettiT_primary_eq3} \\
    a_{00}d_0 + \lambda_1'd_1 &=  \lambda_2'd_2 + \lambda_3'd_3, \label{lemma_bettiT_primary_eq4}
\end{align}
where $ \lambda_i,\lambda_i' < a_{ii}$ for $i = 1,2,3$. Subtracting \eqref{lemma_bettiT_primary_eq3} from \eqref{lemma_bettiT_primary_eq4} gives 
\begin{equation}
    (\lambda_1-\lambda_1')d_1 + (\lambda_2'-\lambda_2)d_2 + (\lambda_3'-\lambda_3)d_3 = 0.
\end{equation}
Since $(-\lambda_1,\lambda_2,\lambda_3) \neq (-\lambda_1',\lambda_2',\lambda_3')$, not all coefficients above are zero. Without loss of generality, suppose $(\lambda_1-\lambda_1') > 0$, $(\lambda_2'-\lambda_2) \leq 0$, $(\lambda_3'-\lambda_3) \leq 0$. This implies $\lambda_1 \geq (\lambda_1-\lambda_1') \geq a_{11}$, which is a contradiction. 
\end{proof}



\begin{lemma}\label{lemma_bettiT_secondary}
     If $S$ is secondary, then $|\mathcal{N}| \leq 1$.
\end{lemma}
\begin{proof} 
By Lemmas \ref{lemma_catenary_lambda1_secondary} and \ref{lemma_catenaryT_secondary_case1}, no points of the form $(-\lambda_1,\lambda_2,\lambda_3)$ or $(\lambda_1,-\lambda_2,\lambda_3)$ (where $\lambda_i \in \mathbb{Z}_{>0}$) are in $\mathcal{N}$.

Thus every point in $\mathcal{N}$ is of the form $(\lambda_1,\lambda_2,-\lambda_3)$ (where $\lambda_i \in \mathbb{Z}_{>0}$). 
If $|\mathcal{N}| \geq 2$, then let $(\lambda_1,\lambda_2,-\lambda_3)$ and $(\lambda_1',\lambda_2',-\lambda_3')$ be distinct points in $\mathcal{N}$. By Lemma \ref{lemma_catenaryT_secondary_case2}, we have
\begin{align}
    a_{00}d_0 + \lambda_3d_3 &=  \lambda_1d_1 + \lambda_2d_2,  \label{lemma_bettiT_secondary_eq1} \\
    a_{00}d_0 + \lambda_3'd_3 &=  \lambda_1'd_1 + \lambda_2'd_2, \label{lemma_bettiT_secondary_eq2}
\end{align}
where $\lambda_3, \lambda_3' < a_{33}$ Since $T = R \cap \{z \leq a_{33}\}$, we have $\lambda_1, \lambda_1' < b_{11}=a_{11}$.  

Without loss of generality, suppose $\lambda_3 \geq \lambda_3'$. Subtracting \eqref{lemma_bettiT_secondary_eq2} and \eqref{lemma_bettiT_secondary_eq1}, we obtain
\begin{equation}
    (\lambda_3- \lambda_3')d_3 + (\lambda_1'-\lambda_1)d_1 =  (\lambda_2-\lambda_2')d_2
\end{equation}
where all above coefficients are nonnegative, as otherwise we would get a contradiction with $\lambda_3 < a_{33}$ or $\lambda_1 < a_{11}$. As $(\lambda_1,\lambda_2,-\lambda_3) \neq (\lambda_1',\lambda_2',-\lambda_3')$, we obtain $\lambda_2 >
(\lambda_2-\lambda_2') \geq a_{22}$. We have 
\begin{equation}
     a_{00}d_0 =  \lambda_1d_1 + (\lambda_2-a_{22})d_2 + (a_{33}-\lambda_3)d_3
\end{equation}
where all above coefficients are positive. This contradicts $\varphi^{-1}(a_{00}d_0) = \{(a_{00},0,0,0), (0,a_{11},0,0)\}$. 
\end{proof}

Let
$$U_0 = \{ \max\{ x,0\} d_1 + \max\{y,0\}d_2 + \max\{z,0\}d_3 \mid (x,y,z) \in V^* \setminus \mathcal{N} \}.$$
Note that $U_0$ is independent of the choice of $V$.

\begin{lemma}\label{lemma_U0_VN}
     $|U_0| = |V^* \setminus \mathcal{N}|$.
\end{lemma}

\begin{proof}
    Suppose that there exist distinct $(-\lambda_1,\lambda_2,\lambda_3)$ and $(-\lambda_1',\lambda_2',\lambda_3') \in V^* \setminus \mathcal{N}$ (where $\lambda_i > 0$) such that $\lambda_2d_2 + \lambda_3d_3 = \lambda_2'd_2 + \lambda_3'd_3$. Without loss of generality, suppose $\lambda_2 > \lambda_2'$ and $\lambda_3 < \lambda_3'$. Now, as $(\lambda_3'-\lambda_3)d_3 = (\lambda_2-\lambda_2')d_2 > 0$, we obtain $\lambda_3' > (\lambda_3'-\lambda_3) \geq a_{33}$. This contradicts $(-\lambda_1',\lambda_2',\lambda_3') \in V^*$, since $T = R \cap \{z \leq a_{33}\}$. 
    
   If there exist distinct $(\lambda_1,-\lambda_2,\lambda_3)$ and $ (\lambda_1',-\lambda_2',\lambda_3') \in V^* \setminus \mathcal{N}$ (where $\lambda_i > 0$) such that $\lambda_1d_1 + \lambda_3d_3 = \lambda_1'd_1 + \lambda_3'd_3$, we proceed analogously as above and derive a contradiction. 
    
   Suppose that there exist distinct $(\lambda_1,\lambda_2,-\lambda_3)$ and $ (\lambda_1',\lambda_2',-\lambda_3') \in V^* \setminus \mathcal{N}$ (where $\lambda_i > 0$) such that $\lambda_1d_1 + \lambda_2d_2 = \lambda_1'd_1 + \lambda_2'd_2$. Without loss of generality, suppose $\lambda_1 > \lambda_1'$ and $\lambda_2 < \lambda_2'$, which gives $\lambda_1 > (\lambda_1 - \lambda_1') \geq a_{11} = b_{11}$ (as $S$ is ordered). This contradicts $(\lambda_1,\lambda_2,-\lambda_3) \in V^*$, since $T \subseteq R \subseteq \{x \leq b_{11}, y \leq b_{22}, z \leq b_{33}\}$.


   Suppose that $S$ is primary and there exist $(-\lambda_1,\lambda_2,\lambda_3)$ and $(\lambda_1',-\lambda_2',\lambda_3') \in V^* \setminus \mathcal{N}$ (where $\lambda_i > 0$) such that $\lambda_2d_2 + \lambda_3d_3 = \lambda_1'd_1 + \lambda_3'd_3$. Without loss of generality, suppose $\lambda_3 \geq \lambda_3'$, which gives $\lambda_1' \geq a_{11}$, which contradicts  $(\lambda_1',-\lambda_2',\lambda_3') \in V^*$, since $T \subseteq \{x \leq a_{11}, y \leq a_{22}, z \leq a_{33}\}$. All other cases follow analogously. 

   Suppose that $S$ is secondary and there exist $(-\lambda_1,\lambda_2,\lambda_3)$ and $(\lambda_1',-\lambda_2',\lambda_3') \in V^* \setminus \mathcal{N}$ (where $\lambda_i > 0$) such that $\lambda_2d_2 + \lambda_3d_3 = \lambda_1'd_1 + \lambda_3'd_3$. Since $(\lambda_1',-\lambda_2',\lambda_3') \in V^*$, we have $\lambda_1' < b_{11} = a_{11}$. Thus we have
   \begin{equation}
       \lambda_2d_2 = \lambda_1'd_1 + (\lambda_3'-\lambda_3)d_3
   \end{equation}
   where $(\lambda_3'-\lambda_3) \geq 0$. This gives $\lambda_2 \geq a_{22}$, hence
   \begin{equation}
       (\lambda_2-a_{22})d_2 = \lambda_1'd_1 + (\lambda_3'-\lambda_3-a_{33})d_3
   \end{equation}
   where $(\lambda_2-a_{22}) \geq 0$. Since $a_{11} > \lambda_1' > 0$, we have $(\lambda_3'-\lambda_3-a_{33}) \geq 0$, which gives $(\lambda_2-a_{22}) \geq a_{22}$. This leads to an infinite descent, which is impossible. 
   All other cases follow analogously.
\end{proof}


Let $\rho$ be a minimal presentation of $S$ and let $\rho^* = \{(a,b) \in \mathbb{N}^{4} \times \mathbb{N}^{4}  \mid (a,b) \in \rho, |\textup{supp}(a) | = | \textup{supp}(b) | =2 \} $. To proceed, we require a result of Bresinsky \cite{Bresinsky1988}. 

\begin{lemma}\label{lemma_rho_bresinsky}
   \textup{\cite[Corollary 2]{Bresinsky1988}} There exists a minimal presentation $\rho$ such that $\eta(S) - |\rho^*| \leq 4$.
\end{lemma}



\begin{lemma}\label{lemma_betti_VNrho_primary}
     If $S$ is primary, then $|U_0| + 4 \geq  \eta(S) \geq |U_0| +3 $. 
\end{lemma}
\begin{proof}
    By Lemma \ref{lemma_catenary_minrelations}, and Theorems \ref{theoremprocedure} and \ref{theorembettifind_primary}, we have
    \begin{equation}\label{lemma_betti_VNrho_primary_eq1}
        \textup{Betti}(S) = U_0 \sqcup \{a_{00}d_0, a_{11}d_1, a_{22}d_2, a_{33}d_3\}. 
    \end{equation}
    Lemma \ref{lemma_catenary_minrelations} tells us that $(a_{00},0,0,0)$ is the only element in its $\mathcal{R}$-class. The same is true for $(0,a_{11},0,0)$, $(0,0,a_{22},0)$, and $(0,0,0,a_{33})$, so the sets on the right-hand side
    of \eqref{lemma_betti_VNrho_primary_eq1} are indeed disjoint. 
    
    If there exists a factorization in $\bigcup_{i=0}^3 \varphi^{-1}(a_{ii}d_i)$ that is distinct from the four canonical ones, then $a_{00}d_0, a_{11}d_1, a_{22}d_2, a_{33}d_3$ together contribute at least three elements to $\rho$ by Lemma \ref{lemma_catenary_minrelations} and Theorem \ref{theorem_minimalpresentation}. Otherwise, since $S$ is not secondary, we have $a_{00}d_0= a_{11}d_1= a_{22}d_2= a_{33}d_3$, which has four $\mathcal{R}$-classes by Lemma \ref{lemma_catenary_minrelations}, so it contributes three elements to $\rho$ by Theorem \ref{theorem_minimalpresentation}.    
    
    
    From Lemmas \ref{lemma_catenary_primary} and \ref{lemma_catenarycase1_primary} we know that every element of $U_0$ has exactly two $\mathcal{R}$-classes, so each contributes exactly one element to $\rho$, by Theorem \ref{theorem_minimalpresentation}. This, combined with the previous paragraph, gives $\eta(S) \geq |U_0| + 3$. 

    Suppose, for contradiction, that $\eta(S) - |U_0| > 4$. By \eqref{lemma_betti_VNrho_primary_eq1}, since $U_0$ contributes exactly $|U_0|$ elements to $\rho$, the elements $a_{00}d_0$, $a_{11}d_1$, $a_{22}d_2$, $a_{33}d_3$ must together contribute more than $4$ elements to $\rho$ --- and by Theorem \ref{theorem_minimalpresentation}, this holds for every minimal presentation $\rho$, since this count is determined only by the number of $\mathcal{R}$-classes. 
    By Lemma \ref{lemma_catenary_minrelations}, there do not exist $z_1,z_2 \in \varphi^{-1}(a_{ii}d_i)$ such that $|\textup{supp}(z_1)| = |\textup{supp}(z_2)| = 2$ and $\textup{supp}(z_1) \cap \textup{supp}(z_2) = \varnothing$. Thus in every minimal presentation $\rho$, there will be more than $4$ elements that are not in $\rho^*$. This gives a contradiction with Lemma \ref{lemma_rho_bresinsky}, therefore $\eta(S) \leq |U_0| + 4 $.
\end{proof}

\begin{lemma}\label{lemma_betti_VNrho_secondary}
    If $S$ is secondary, then $|U_0| + 6 \geq  \eta(S) \geq |U_0| + 2 $. 
\end{lemma}
\begin{proof}
    By Lemma \ref{lemma_catenary_minrelations} and Theorems \ref{theoremprocedure} and \ref{theorembettifind_secondary}, we have
   \begin{equation}\label{lemma_betti_VNrho_secondary_eq1}
        U_0 \sqcup \{a_{00}d_0, a_{22}d_2\} \subseteq \textup{Betti}(S) \subseteq U_0 \cup \{a_{00}d_0, a_{22}d_2, b_{22}d_2, b_{33}d_3\}. 
    \end{equation}
     From Lemmas \ref{lemma_catenary_lambda1_secondary} and \ref{lemma_catenary_lambda23_secondary} we know that every element of $U_0$ has exactly two $\mathcal{R}$-classes, so it contributes exactly one element to $\rho$, by Theorem \ref{theorem_minimalpresentation}. We also have $((a_{00},0,0,0),(0,a_{11},0,0) ) \in \rho$ and $((0,0,a_{22},0),(0,0,0,a_{33})) \in  \rho$, hence $\eta(S) \geq |U_0| + 2$.

     We claim that $b_{22}d_2$ and $b_{33}d_3$ together contribute at most four elements to $\rho$. Otherwise, without loss of generality, suppose that $b_{22}d_2$ contributes at least three elements to $\rho$, so it has at least four $\mathcal{R}$-classes. The only way for this to be possible is $\varphi^{-1}(b_{22}d_2) = \{(b_{20},0,0,0)$, $(0,b_{21},0,0)$, $(0,0,b_{22},0)$, $(0,0,0,b_{23})\}$. Note that $b_{23} \geq b_{33}$. However, if $b_{23} > b_{33}$, then $(b_{30}$, $b_{31}$, $b_{32}$, $(b_{23}-b_{33})) \in \varphi^{-1}(b_{22}d_2)$, which is a contradiction, since $b_{30}, (b_{23}-b_{33}) > 0 $. Hence $b_{23}=b_{33}$ and $b_{22}d_2 = b_{33}d_3$. Therefore, in the current case, $b_{22}d_2$ and $b_{33}d_3$ together contribute at most three elements to $\rho$, by Theorem \ref{theorem_minimalpresentation}, since $b_{22}d_2 = b_{33}d_3$ has at most four $\mathcal{R}$-classes. 
     
     
    Therefore, by \eqref{lemma_betti_VNrho_secondary_eq1}, we have $\eta(S) \leq |U_0| + 2 + 4 = |U_0| + 6$.   
\end{proof}







\begin{cor}\label{corollary_type_main}
     If $S$ is primary, then $4 t(S) +5 \geq \eta(S) \geq t(S) - 11$.  If $S$ is secondary, then $2 t(S) +9 \geq \eta(S) \geq t(S) - 4$. 
\end{cor}

\begin{proof}
    Suppose that $S$ is primary. The lower bound for $\eta(S)$ follows from the chain of inequalities:
    \begin{equation}
        \eta(S) \geq |U_0| + 3 = |V^* \setminus \mathcal{N}| + 3 \geq |V^*|  \geq t(S) - 11
    \end{equation}
     which are the result of Lemmas \ref{lemma_betti_VNrho_primary}, \ref{lemma_U0_VN}, \ref{lemma_bettiT_primary}, and Theorem \ref{theorem_type1}, respectively. Moreover, the upper bound for $\eta(S)$ is due to the following chain of inequalities:
    \begin{equation}
        \eta(S) \leq |U_0| + 4 = |V^* \setminus \mathcal{N}| + 4 \leq |V^*| +4 \leq 4t(S) +5
    \end{equation}
     which are the result of Lemmas \ref{lemma_betti_VNrho_primary}, \ref{lemma_U0_VN}, and Theorem \ref{theorem_type1}, respectively.

     Suppose that $S$ is secondary. The lower bound for $\eta(S)$ is due to the following chain of inequalities:
    \begin{equation}
        \eta(S) \geq |U_0| + 2 = |V^* \setminus \mathcal{N}| + 2 \geq |V^*| +1 \geq t(S) - 4
    \end{equation}
     which are the result of Lemmas \ref{lemma_betti_VNrho_secondary}, \ref{lemma_U0_VN}, \ref{lemma_bettiT_secondary}, and Theorem \ref{theorem_type1}, respectively. Moreover, the upper bound for $\eta(S)$ is due to the following chain of inequalities:
    \begin{equation}
        \eta(S) \leq |U_0| + 6 = |V^* \setminus \mathcal{N}| + 6 \leq |V^*| +6 \leq 2t(S) +9
    \end{equation}
     which are the result of Lemmas \ref{lemma_betti_VNrho_secondary}, \ref{lemma_U0_VN}, and Theorem \ref{theorem_type1}, respectively. 
\end{proof}


For $t(S) \geq 2$, the inequality $4t(S) +5 \geq 2 t(S) + 9$ holds, so Corollary \ref{corollary_type_main} implies Theorem \ref{theorem_type_main} in this case. Recall that $t(S) = 1$ is equivalent to $S$ being symmetric \cite{Kunz1970}. This allows us to use the following result to resolve the case $t(S) = 1$. 


\begin{thm}
    \textup{\cite{Bresinsky1975sym}} If $S$ is symmetric with $e(S)=4$, then $\eta(S) \in \{3,5\}$.
\end{thm}

Since $5 \leq 9 = 4 \cdot 1 + 5$, this completes the proof of Theorem \ref{theorem_type_main}.


\section{Closing remarks}
The bounds of Theorem \ref{theorem_type_main} can most likely be improved. To improve the bound $\eta(S) \geq t(S)-11$, we propose the following conjecture.


\begin{conj}\label{conjecture}
    If $S$ is a numerical semigroup with $e(S) = 4$, then $\eta(S) \geq t(S) +1$.
\end{conj}
 
We expect that $\eta(S) \geq t(S) +1$ is the best possible bound, since no numerical semigroup with $e(S)=4$ and $\eta(S) \leq t(S)$ is currently known. In contrast, the situation where $\eta(S) = t(S) +1$ is quite common among numerical semigroups with embedding dimension four. For example, $S = \langle 77,85,89,131 \rangle$ has $e(S)=4$, $t(S)=3$, and $\eta(S)=4$. Likewise, $S = \langle 76,99,143,223 \rangle$ has $e(S)=4$, $t(S)=6$, $\eta(S)=7$, and $S = \langle 195,200,209,226 \rangle$ has $e(S)=4$, $t(S)=8$, $\eta(S)=9$. The equation $\eta(S) = t(S) +1$ is also satisfied for all numerical semigroups with $e(S)=3$ (Theorem \ref{theorem_embd3}).

\begin{remark}
    In higher embedding dimension, the inequality $\eta(S) \geq t(S) +1$ is false in general, already for embedding dimension five. For example, 
    $S=\langle 86,133,143,155,198 \rangle$ has $e(S)=5$ and $t(S)=\eta(S)=10$. Likewise, $S=\langle 161,177,196,207,225 \rangle$ has $e(S)=5$, $t(S)= 18$, $\eta(S)=16$, and (a more extreme example) $S = \langle 6072, 6307,7965, 8702, 9423 \rangle$ has $e(S)=5$, $t(S)=53$, $\eta(S)=40$.
\end{remark}

\begin{remark}
    The results of Section \ref{section_geometricprocedure}, together with the results of Section \ref{section_minimal_presentations} concerning primary semigroups, appeared earlier in the preprint \cite{Chomicz}.
\end{remark}

\section{Acknowledgments}
At the time of writing this paper, the author was a first-year student at Jagiellonian University in Kraków. The author thanks his tutor,
Tomasz Kowalczyk, for his support.
The author also thanks Alessio Moscariello and Alessio Sammartano for their comments.








\bibliographystyle{plain}
\bibliography{refs.bib}

\vspace{5pt}

%\noindent
%\textbf{Statements and Declarations}
%\\ \\
%\textbf{Funding} The author declares that no funds, grants, or other support were received %during the preparation of this manuscript. 
%\\ \\
%\textbf{Competing Interests} The author has no relevant financial or non-financial interests to disclose.
%\\ \\
%\textbf{Data Availability Statement} The author declares that the data supporting the findings of this study are available within the paper. The parts calculated using the algebraic calculator Macaulay2 can be replicated from the data provided. At the end of this article an example code from Macaulay2 is provided.


%\vspace{10pt}

\begin{small}





\noindent
Kazimierz Chomicz

\noindent
Institute of Mathematics

\noindent
Faculty of Mathematics and Computer Science

\noindent
Jagiellonian University

\noindent
ul. Łojasiewicza 6, 30-348 Kraków, Poland

\noindent
e-mail: kazimierz.chomicz@student.uj.edu.pl



\end{small}	










  
\end{document}


